\documentclass[11pt]{article}
\usepackage[a4paper,margin=24mm]{geometry}
\usepackage[T1]{fontenc}
\usepackage{amsmath,amsfonts,amssymb,amsthm,lmodern,microtype}
\usepackage{booktabs,tabularx,array,graphicx,xcolor,verbatim}
\usepackage{algorithm,algpseudocode,enumitem}
\usepackage[font=small,labelfont=bf]{caption}
\usepackage[hidelinks]{hyperref}
\usepackage{fancyhdr}
\pagestyle{fancy}\fancyhf{}\fancyhead[L]{\small Verified Context Retirement}\fancyhead[R]{\small Research manuscript v8.0}\fancyfoot[C]{\thepage}
\setlength{\headheight}{14pt}\setlength{\parskip}{3pt}\setlength{\emergencystretch}{2em}
\newtheorem{theorem}{Theorem}\newtheorem{proposition}[theorem]{Proposition}\newtheorem{lemma}[theorem]{Lemma}
\theoremstyle{definition}\newtheorem{definition}[theorem]{Definition}
\theoremstyle{remark}\newtheorem{remark}[theorem]{Remark}
\newcommand{\PASS}{\ensuremath{\mathsf{PASS}}}\newcommand{\FAIL}{\ensuremath{\mathsf{FAIL}}}\newcommand{\UNKNOWN}{\ensuremath{\mathsf{UNKNOWN}}}
\newcommand{\T}{\mathcal T}
\newcommand{\E}{\mathbb E}\newcommand{\Prb}{\mathbb P}\newcommand{\ind}{\mathbf 1}\newcommand{\R}{\mathcal R}
\DeclareMathOperator*{\argmin}{arg\,min}
\hypersetup{pdftitle={Verified Context Retirement: Safety Conditions and Retrieval-Memory Trade-offs},pdfauthor={Research manuscript; author information pending confirmation}}
\title{\vspace{-9mm}\textbf{Verified Context Retirement:\\[1mm] Safety Conditions and Retrieval--Memory Trade-offs}}
\author{Research manuscript\\[1mm]
\small This artifact candidacy draft accompanies the \texttt{future-code} research repository. Author names and affiliations require confirmation before any submission. Mathematical development, implementation, experiments, and writing were AI-assisted. This is a research artifact, not an accepted or peer-reviewed venue.}
\date{September 28, 2026 \quad | \quad Version 8.0 (major revision: scoped claims, consistency audit)}
% Auto-generated from the frozen cohort (do not edit).
\newcommand{\RetNullClosure}{\textbf{triggered}}
\newcommand{\RetAttributionClosure}{\textbf{triggered}}
\newcommand{\RetHactClosure}{\textbf{not triggered}}
\newcommand{\RetFrontierVsFull}{58.8\%}
\newcommand{\RetFrontierVsMaskingBytes}{-4.4\%}
\newcommand{\RetOpsReduction}{48.3\%}
\newcommand{\RetMakespanReduction}{87.4\%}
\newcommand{\RetEvidenceReduction}{59.8\%}
\newcommand{\RetSummaryResetAccepted}{21/63}
\newcommand{\RetSummaryResetStale}{30{,}820}
\newcommand{\RetSummaryResetLost}{280}
\newcommand{\RetFrontierAccepted}{63/63}

% Auto-generated from the frozen v2 cohort (do not edit).
\newcommand{\RetwoNullClosure}{\textbf{not triggered}}
\newcommand{\RetwoAttributionClosure}{\textbf{not triggered}}
\newcommand{\RetwoHactClosure}{\textbf{pass}}
\newcommand{\RetwoAdaptivePromotable}{\textbf{not promotable}}
\newcommand{\RetwoAuditClosure}{\textbf{clean}}
\newcommand{\RetwoNullReduction}{29.6\%}
\newcommand{\RetwoNullThreshold}{20\%}
\newcommand{\RetwoBytesDelta}{-8.0\%}
\newcommand{\RetwoOpsReduction}{44.0\%}
\newcommand{\RetwoCostOReduction}{29.6\%}
\newcommand{\RetwoMakespanReduction}{87.5\%}
\newcommand{\RetwoRuns}{1176}
\newcommand{\RetwoCells}{84}
\newcommand{\RetwoLongRuns}{63}
\newcommand{\RetwoLongOpsReduction}{45.1\%}
\newcommand{\RetwoLongCostOReduction}{30.9\%}
\newcommand{\RetwoSensK}{17.8\%}
\newcommand{\RetwoSensM}{29.6\%}
\newcommand{\RetwoSensXxl}{38.8\%}
\newcommand{\RetwoSummaryResetAccepted}{28/84}
\newcommand{\RetwoSEffectOnBytes}{0.000\%}
\newcommand{\RetwoDensityRegret}{0.0\%}
\newcommand{\RetwoAdaptiveReduction}{-0.0\%}
\newcommand{\RetwoWorstFamily}{none}
\newcommand{\RetwoFrontierAccepted}{84/84}
\newcommand{\RetwoBytesWin}{8.0\%}
\newcommand{\RetwoCOEightPooled}{1{,}364\,B}
\newcommand{\RetwoCOEightFivePooled}{658\,B}
\newcommand{\RetwoCOEightClustered}{1{,}047\,B}
\newcommand{\RetwoCOEightIndependent}{191\,B}
\newcommand{\RetwoCOEightGlobal}{never}
\newcommand{\RetwoGlobalRatio}{0.924}

% Auto-generated from the frozen v2 cohort via the pre-specified
% exploratory analyses E1, E2, E3, E6 (docs/protocols/EXPLORATORY_ANALYSES_v2.md).
% These are exploratory (post-data, pre-specified), not confirmatory.
\newcommand{\ExpFrontierMaskingSignO}{84/84 cells ($p=1.0\times10^{-25}$)}
\newcommand{\ExpFrontierMaskingSignB}{76/84 cells ($p=5.0\times10^{-15}$)}
\newcommand{\ExpFrontierBhSignO}{34/84 cells ($p=0.10$)}
\newcommand{\ExpFrontierBhSignB}{28/84 cells ($p=3.0\times10^{-3}$)}
\newcommand{\ExpMaskingBhSignO}{0/84 cells ($p=1.0\times10^{-25}$)}
\newcommand{\ExpSeedOpsMin}{33.8\%}
\newcommand{\ExpSeedOpsMax}{50.6\%}
\newcommand{\ExpSeedOMin}{23.3\%}
\newcommand{\ExpSeedOMax}{33.7\%}
\newcommand{\ExpBhDueExposed}{69.0\%}
\newcommand{\ExpBhFootprintExposed}{35.9\%}
\newcommand{\ExpBhBreakEvenPooled}{24{,}444\,B}
\newcommand{\ExpBhBreakEvenIndependent}{2{,}148\,B}
\newcommand{\ExpBhBreakEvenGlobal}{never}
\newcommand{\ExpTraceCommits}{48}
\newcommand{\ExpTraceRegistry}{2{,}077}
\newcommand{\ExpTraceDensity}{0.042}
\newcommand{\ExpTraceHactRatio}{1.71}
\newcommand{\ExpTraceRuleRegret}{matched}
\newcommand{\ExpTraceOpsCut}{23.0\%}
\newcommand{\ExpTraceOpsRange}{18.9\%--25.9\%}
\newcommand{\ExpTraceOCut}{19.6\%}
\newcommand{\ExpTraceAcceptance}{7/7 seeds, both arms}
\newcommand{\ExpLhBhClusteredO}{3{,}355{,}230}
\newcommand{\ExpLhBhIndependentO}{44{,}428{,}129}
\newcommand{\ExpLhBhGlobalO}{2{,}097{,}520}
\newcommand{\ExpLhBhClusteredB}{2{,}170{,}316}
\newcommand{\ExpLhBhIndependentB}{4{,}182{,}003}
\newcommand{\ExpLhBhGlobalB}{2{,}097{,}520}
\newcommand{\ExpLhBhClusteredOps}{289}
\newcommand{\ExpLhBhIndependentOps}{9{,}826}
\newcommand{\ExpLhBhGlobalOps}{0}

\begin{document}
\maketitle
\begin{abstract}
When can an agent stop remembering its execution history without
losing the current evidence and obligations it needs to finish
correctly---and in which cost regime might the answer pay? We study
this question at the \emph{mechanism level}: a pre-registered,
factorial, frozen-protocol simulator study of context retirement,
with 1{,}176 runs crossing seven retention policies with a scheduler
factor, plus a 160-episode long-horizon sub-cohort, across four
registry sizes, three dependency families, and seven seeds, with
temporal persistence, obligation lifetimes, three-phase accounting, a
five-episode capsule audit against stream ground truth, and two frozen
cost regimes (bytes; bytes plus a per-round-trip constant). Within
this simulator, the pattern that survives is a \emph{conversion
heuristic with a regime-conditional payoff}. Every measured policy
that retires history through a lossy summary fails the acceptance
contract (periodic summary/reset accepts 28 of 84 runs, with
109{,}761 stale reads and 385 lost obligations), while the policy that
first converts history into verified interface bindings, a typed
pending-obligation set, and declared dependencies accepts 84/84 and
passes the capsule audit with zero defective service. The payoff
direction in the simulator is operations, not bytes: frontier
retirement cuts on-demand retrieval operations by \RetwoOpsReduction{}
pooled ($\RetwoLongOpsReduction{}$ at four times the horizon; sign
test $p{=}10^{-25}$ on round-trip-regime cost) and is cheaper than
masking on bytes by \RetwoBytesWin{} pooled---below the frozen 20\%
byte gate, so no byte-threshold win is claimed. A closed-form
break-even rule (Proposition~\ref{prop:breakeven}) converts the frozen
sensitivity sweep into a conditional threshold formula
($c_{0.8}\approx1{,}364$\,B pooled in the simulator); it is a property
of the frozen accounting, not validated deployment guidance.
Sufficiency of the conversion heuristic is proved \emph{within the
model} (Proposition~\ref{prop:sufficiency}): under the engine's stream
invariants and reliable retrieval, the frontier policy provably
incurs zero staleness and zero obligation loss---a structural
prediction of the simulator, not a guarantee about agents. The
scheduler factor is byte-neutral by design and measured at
\RetwoSEffectOnBytes{} leak with an \RetwoMakespanReduction{} makespan
reduction. Grouped certificate layouts survive only as a conditional
component selected by measured update density (dense stratum pass;
\RetwoDensityRegret{} decision-rule regret), and a demand-driven
adaptive-coordination variant was evaluated against pre-registered
promotion gates and \emph{rejected} ($\RetwoAdaptiveReduction{}$
versus frontier). Pre-specified exploratory analyses (labeled as
such) stress the result: a bounded-recency comparator is cheaper than
frontier below a stated break-even price but leaves 69.0\% of due
obligations outside its retention, quantifying what it delegates to
the retrieval layer; and a replay of a real repository's change
history through the same engine preserves the operations advantage
(23.0\% cut) with the density rule matching the oracle on its first
real trace. All experiments are simulator-level: no hosted model was
called, no token or billing claim is made, and the live-agent
cohort---which measures retrieval, context, acceptance, and
end-to-end outcomes under matched budgets---remains future work with
its protocol pre-registered here.
\end{abstract}

\noindent\textbf{Keywords:} bounded evidence; verification contracts; bandwidth-limited logging; completion hyperedges; cost-aware layout selection.

\section{Productivity Boundary and Research Scope}\label{sec:productivity}
The governing research objective is to reduce the cost and time
required for agents to complete large software projects correctly,
without increasing coordination failures or weakening acceptance
contracts. That objective outruns what any single paper can measure,
and this manuscript is explicit about the ladder of evidence it
occupies.

\paragraph{What this paper is.}
It is a mechanism study of context retirement
(Section~\ref{sec:retire}) plus the verification and evidence
components the surviving rule depends on
(Section~\ref{sec:substrate}, with derivations in the appendices). Two protocol generations
were run; the current one is factorial, pre-registered, and
git-frozen before data generation. Five closure rules were declared
before data; two of the paper's own proposed extensions failed their
gates and are reported as rejected: the byte-regime productivity
claim for frontier retirement (it saves \RetwoBytesDelta{} on
submitted bytes versus masking---no win claimed) and the demand-driven
adaptive-coordination variant ($\RetwoAdaptiveReduction{}$ regime-O
versus plain frontier; promotion gate $\ge$10\%). What survives is
(i)~a necessary conversion rule---retire history only after
converting it into verified bindings, typed obligations, and declared
dependencies; lossy conversion fails the acceptance contract at every
scale tested (28/84 at the largest grid); (ii)~an operations result
at equal correctness---frontier retirement cuts on-demand retrieval
operations \RetwoOpsReduction{} and round-trip-regime cost
\RetwoCostOReduction{} versus observation masking, persisting at
$4\times$ horizon, with a clean capsule audit; and (iii)~a
density-selected conditional role for certificate-tree evidence
layouts (\RetwoDensityRegret{} decision-rule regret). HACT is not the
paper's center; it survives only where its ablation says it wins.

\paragraph{The boundary that remains.}
The retirement study is mechanism-level. It models information flow
between a context policy and a seeded work stream; it does not call a
hosted model, and it cannot establish that any agent finishes faster,
bills less, or reasons better. Its staleness, obligation-loss, and
audit results are structural predictions about policies, not
measurements of agents. The live-agent cohort that could test those
predictions has its protocol pre-registered in
Section~\ref{sec:live-protocol} and has not been run. The negative
controls below---flat checker timings, per-episode verification
overhead, unchanged root-status inputs---are reported as limits of
the serialization claim, and the retirement cohort's own closures are
reported as limits of the mechanism claim. Nothing in this paper is
an agent-productivity measurement, and the abstract says so.

\section{Introduction}\label{sec:intro}
Long-horizon software agents accumulate execution history, and the
accumulation is not free: resubmitting it costs bytes that grow with
project length, and naively discarding it costs correctness. The
research question of this paper is the review's stronger form:

\begin{quote}
When can an agent stop remembering its execution history without
losing the current evidence and obligations needed to finish
correctly---and in which cost regime does the answer pay?
\end{quote}

We answer it with a pre-registered, factorial mechanism study
(Section~\ref{sec:retire}). Seven retention policies---from full
replay through observation masking, periodic summary/reset, bounded
pages, and a verified board, to verified-frontier retirement with
typed capsules, plus a demand-driven adaptive variant---are crossed
with a scheduler factor in a full factorial, and consume
byte-identical seeded work streams with temporal persistence,
payload-bearing facts, persistent dependency footprints, and
obligations that come due later. Accounting is split into
construction, steady state, and a five-episode audit tail that
compares each policy's capsule against stream ground truth; two cost
regimes are frozen ex ante, one byte-bound and one round-trip-bound.
The v2 protocol and its five closure rules were committed to git
before any v2 data was generated; the cohort and its tables were
committed before any v2 manuscript text was written. Two protocol
generations were run; the first's design limitations are stated in
Section~\ref{sec:retire-provenance} and its full record is retained in
Appendix~\ref{app:retire-v1}.

The pattern that survives both generations in this simulator is a
conversion heuristic with a
regime-conditional payoff. \emph{Failure direction (measured)}:
every policy that
converts lossily---periodic summary/reset---fails the frozen acceptance
contract at every scale measured (28/84 runs at the v2 scale; 21/63 at
the v1 scale), silently reverting values to beliefs older than the
working window and dropping obligations that cross reset boundaries;
every measured policy that first converts history into verified
interface bindings, a typed pending-obligation set, and declared
dependencies passes it.
\emph{Payoff}: the payoff direction is operations, not bytes---but the byte
direction is also positive. Frontier retirement is cheaper than
masking on bytes by \RetwoBytesWin{} pooled (76 of 84 cells,
$p{=}5{\times}10^{-15}$; below the frozen 20\% byte gate) and on
round-trip-regime cost by \RetwoCostOReduction{} (84 of 84 cells,
$p{=}10^{-25}$), cutting operations \RetwoOpsReduction{} versus
masking with the advantage persisting at four times the horizon
(\RetwoLongOpsReduction{}). A closed-form break-even rule
(Proposition~\ref{prop:breakeven}) converts the frozen sensitivity
sweep into a conditional threshold formula: the simulator's gate
flips at round-trip price $c_{0.8}\approx1{,}364$\,B pooled on the
frozen cohort---a property of the frozen accounting, not validated
deployment guidance. The
live-agent hypothesis this scopes is precise: savings, if any,
come from avoided tool-call round trips; the byte direction is a
bonus too small to carry the claim. \emph{Audit}: the promoted
mechanism passes a capsule audit against stream ground truth with zero
defective service, and the audit's development history includes a
disclosed, test-caught reference fix (Section~\ref{sec:retire-results}).
\emph{Sufficiency (within the model)}: the conversion heuristic is
provably sufficient under the engine's semantics
(Proposition~\ref{prop:sufficiency}): given the stream invariants and
reliable retrieval, the frontier policy incurs zero staleness and
zero obligation loss. This is a model-internal guarantee---a
structural prediction of the simulator---not a claim about agents.
\emph{Stress}: pre-specified exploratory analyses
(Section~\ref{sec:retire-exploratory}) put the result under load: a
bounded-recency window beats frontier below a break-even price of
\ExpBhBreakEvenPooled{} but leaves \ExpBhDueExposed{} of due
obligations outside its retention; a real repository trace replayed
through the engine preserves the operations cut (\ExpTraceOpsCut{})
with the density rule matching the oracle; and the per-seed range of
the operations cut never crosses zero.
\emph{Rejected extensions}: a demand-driven adaptive-coordination
variant was evaluated against a pre-registered promotion gate and
rejected ($\RetwoAdaptiveReduction{}$ regime-O versus plain frontier;
the gate required $\ge$10\%), and the HACT certificate-tree evidence
layout remains a conditional component---against a steel-manned flat
ledger it wins dense updates and loses sparse ones, so it is selected
by a measured-density rule with \RetwoDensityRegret{} regret rather
than by default. The scheduler factor is cleanly separated: measured
byte leak \RetwoSEffectOnBytes{}, makespan reduction
\RetwoMakespanReduction{}.

Section~\ref{sec:substrate} condenses the verification and
evidence components the conversion heuristic relies on, unchanged in
substance: immutable identities, packet cost models, layout learning,
and the negative controls that bound what serialized evidence can
claim; its full derivations, including the exact-DP reference
formulation, now live in the appendices. A capsule that carries verified bindings and pending
obligations needs an evidence layout when its contents must cross a
link, and the layout's value is conditional in exactly the way the
ablation measures. Sections~\ref{sec:delivery}--\ref{sec:discussion}
bound the delivery, adoption, and trust claims.

The negative controls remain: flat checker timings and per-episode
verification overhead are limits of the serialization claim; the
mechanism-level cohort is a limit of the retirement claim; the
live-agent protocol of Section~\ref{sec:live-protocol} has not been
run, and no hosted model, token count, or billing figure appears
anywhere in this paper as a measurement. Executing that
pre-registered protocol is the study's highest-priority next
experiment.

\section{Related Work and Positioning}\label{sec:related}
\paragraph{Context management for long-horizon agents.}
Observation masking with focused retrieval is the strongest simple
baseline for bounded context~\cite{complexitytrap}; context
folding~\cite{contextfolding} branches and folds trajectories; periodic
summary/reset is the generic compaction baseline. Our mechanism study
steel-mans all three in a factorial design and reports where each
fails: masking is near-optimal on submitted bytes but pays for every
non-current read
with a retrieval operation (\RetwoOpsReduction{} more operations than
frontier retirement at equal correctness), while summary/reset
---cheapest on bytes---fails the acceptance contract on 56 of 84 runs
through silent staleness and lost obligations. This locates our contribution
precisely: not a better compaction heuristic, but a measured
characterization of \emph{which retained structures} (typed pending
obligations, current verified bindings, dependency interfaces) make
forgetting safe, and a pre-registered operations metric for the live
cohort that could confirm it.

\paragraph{Multi-agent coordination.}
Decentralized shared context~\cite{delm} and adaptive communication
topologies address central-manager bottlenecks; controlled
scaling studies~\cite{scale} show coordination can hurt sequential or
tool-heavy work. Our verified-board arm is exactly such a baseline and
holds correctness (84/84) at the cost of shipping a full present-state
row per gate; the scheduler factor is orthogonal and byte-neutral by
design (measured leak \RetwoSEffectOnBytes{}), consistent with the
task-architecture alignment those studies report.

\paragraph{Testing, impact analysis, and incremental computation.}
Test prioritization seeks earlier fault exposure~\cite{rothermel};
regression selection uses change dependencies~\cite{ekstazi,namerts}.
Our dependency footprints and completion records inherit that lineage;
the certificate tree is an incremental-publication structure, not a
selection structure, and its new ablation shows its value is
conditional on update density (dense global updates favor the tree by
$91.3\%$ at the v2 scale; sparse deltas favor a flat ledger).

\paragraph{Benchmarks.}
SWE-bench~\cite{swebench} and its harness~\cite{swebenchharness}
provide the external anchor the live cohort pre-registers in
Section~\ref{sec:live-protocol}; no benchmark result appears in this
paper because no live agent was run.

\paragraph{What is not claimed.}
Serialization savings are not compression savings: fewer padded bytes
need not imply fewer compressed bytes, and we measure both where both
exist. Flat ledgers and a trusted root-status endpoint can be much
smaller than an archive of aggregation updates; they satisfy different
delivery needs and are steel-manned, not hidden. No claim in this
paper extends to agent speed, billing, or reasoning quality; those
belong to the unrun live cohort, whose protocol is pre-registered with
closure rules.

% Auto-generated from the frozen v2 cohort (do not edit).
% Columns: policy, accepted, regime-B bytes/run, regime-O cost/run, ops/run,
% A3 audit defects, staleness defects, obligation losses.
\newcommand{\RetwofullreplayRow}{full-replay & 84/84 & 12{,}007{,}957 & 12{,}007{,}957 & 0 & 0 & 0 & 0 \\}
\newcommand{\RetwomaskingRow}{masking & 84/84 & 5{,}107{,}068 & 12{,}786{,}190 & 1{,}875 & 0 & 0 & 0 \\}
\newcommand{\RetwosummaryresetRow}{summary-reset & 28/84 & 220{,}410 & 220{,}410 & 0 & 0 & 109{,}761 & 385 \\}
\newcommand{\RetwoboundedhierRow}{bounded-hier & 84/84 & 789{,}534 & 5{,}748{,}522 & 1{,}211 & 0 & 0 & 0 \\}
\newcommand{\RetwoverifiedboardRow}{verified-board & 84/84 & 12{,}929{,}557 & 12{,}929{,}557 & 0 & 0 & 0 & 0 \\}
\newcommand{\RetwofrontierRow}{frontier & 84/84 & 4{,}698{,}273 & 9{,}002{,}292 & 1{,}051 & 0 & 0 & 0 \\}
\newcommand{\RetwofrontieradaptiveRow}{frontier-adaptive & 84/84 & 4{,}698{,}350 & 9{,}005{,}733 & 1{,}052 & 0 & 0 & 0 \\}

% Auto-generated from the frozen v2 cohort (do not edit).
\newcommand{\RetwoStratumClusteredRow}{clustered & 705{,}448 & 676{,}449 & 39.8\% & 230 & 77 \\}
\newcommand{\RetwoStratumIndependentRow}{independent & 2{,}607{,}799 & 2{,}323{,}054 & 39.6\% & 5{,}394 & 3{,}075 \\}
\newcommand{\RetwoStratumGlobalRow}{global & 12{,}007{,}957 & 11{,}095{,}317 & 7.6\% & 0 & 0 \\}

% Auto-generated from the frozen v2 cohort (do not edit).
\newcommand{\RetwoLongClusteredmaskingRow}{clustered & masking & 7/7 & 745 & 4{,}449{,}648 & 1{,}398{,}713 \\}
\newcommand{\RetwoLongClusteredfrontierRow}{clustered & frontier & 7/7 & 284 & 2{,}559{,}028 & 1{,}394{,}594 \\}
\newcommand{\RetwoLongClusteredfrontieradaptiveRow}{clustered & frontier-adaptive & 7/7 & 287 & 2{,}569{,}097 & 1{,}394{,}715 \\}
\newcommand{\RetwoLongIndependentmaskingRow}{independent & masking & 7/7 & 16{,}051 & 72{,}702{,}403 & 6{,}956{,}921 \\}
\newcommand{\RetwoLongIndependentfrontierRow}{independent & frontier & 7/7 & 8{,}943 & 42{,}826{,}402 & 6{,}194{,}704 \\}
\newcommand{\RetwoLongIndependentfrontieradaptiveRow}{independent & frontier-adaptive & 7/7 & 8{,}958 & 42{,}888{,}791 & 6{,}196{,}238 \\}
\newcommand{\RetwoLongGlobalmaskingRow}{global & masking & 7/7 & 0 & 34{,}003{,}520 & 34{,}003{,}520 \\}
\newcommand{\RetwoLongGlobalfrontierRow}{global & frontier & 7/7 & 0 & 31{,}417{,}920 & 31{,}417{,}920 \\}
\newcommand{\RetwoLongGlobalfrontieradaptiveRow}{global & frontier-adaptive & 7/7 & 0 & 31{,}417{,}920 & 31{,}417{,}920 \\}

% Auto-generated from the frozen v2 cohort (do not edit).
\newcommand{\RetwoEvClusteredRow}{clustered & 150{,}386 & 204{,}631 & 36.1\% \\}
\newcommand{\RetwoEvIndependentRow}{independent & 277{,}129 & 344{,}823 & 24.4\% \\}
\newcommand{\RetwoEvGlobalRow}{global & 1{,}936{,}704 & 169{,}408 & -91.3\% \\}

\section{When Can an Agent Stop Remembering? A Pre-Registered Mechanism Study}\label{sec:retire}

\subsection{The question}\label{sec:retire-question}
Long-horizon agents face a retention decision every episode: resubmit
history, mask it and re-fetch on demand, or summarize it away. The
research question is when the third option---\emph{retirement}---is
safe, and what it buys:
\begin{quote}
\emph{When can an agent stop remembering its execution history without
losing the current evidence and obligations it needs to finish
correctly, and in which cost regime does the answer pay?}
\end{quote}

The study is \emph{mechanism-level}: every number below is produced by
a deterministic engine that models information flow between a context
policy and a seeded work stream. No hosted model is called. The
answer that survives two frozen protocol generations is a
\textbf{conversion rule with a regime-conditional payoff}:
\begin{quote}
\emph{An agent may retire its history only after converting it into
(i)~verified interface bindings for every fact another component will
read, (ii)~a typed pending-obligation set, and (iii)~declared
dependency interfaces. The payoff is not bytes: it is avoided
retrieval round trips, which pay only in the round-trip-bound regime.}
\end{quote}

\subsection{Study provenance and what changed between generations}\label{sec:retire-provenance}
Two protocol generations were run, both pre-registered.
\paragraph{Generation v1 (file-frozen 2026-09-23).} Ten hand-picked
arms over $3\times3\times7=63$ cells (630 runs) with three closure
rules. Its closures fired: no byte-cost win for frontier retirement
over masking (null closure), an attribution closure that forced the
dimension-separated reading, and a conditional pass for grouped
certificate layouts. v1's protocol and engine were frozen in files
before its cohort ran, but were committed to git only on 2026-09-24,
after the v2 protocol commit; v1's freeze is therefore file-attested
only, and this ordering is disclosed here rather than smoothed over.
The full v1 design, its numbers, and its deviations remain in
Appendix~\ref{app:retire-v1}.
\paragraph{Generation v2 (git-frozen 2026-09-24, commit before data).}
v1 answered ``what does each policy lose'' but had four design
weaknesses: a hand-picked arm list rather than a factorial, a single
pooled byte threshold, capsule accounting evaluated only as an
end-of-run total, and one cost regime. v2 repairs all four ex ante:
(1)~a full factorial over retention $R$ (7 levels), scheduler $S$
(2 levels), with factor contrasts replacing pairwise arm comparisons;
(2)~three-evaluation separation---every run's accounting is split into
construction (episodes 0--4), steady state (5--34), and a five-episode
audit tail (35--39) in which an independent auditor compares the
policy's capsule against stream ground truth (A1 stale retention, A2
pending-set mismatch, A3 defective service); (3)~two frozen cost
regimes---bandwidth-bound $B$ (bytes) and round-trip-bound $O$ (bytes
$+4{,}096$\,B per retrieval operation, with sensitivity at 1{,}024 and
16{,}384\,B); (4)~scale: four registry sizes ($n\in\{128,512,1{,}024,
4{,}096\}$), three families, seven seeds, forty episodes
($4\times3\times7=84$ cells $\times\,14$ policy runs $=1{,}176$ runs),
plus a long-horizon sub-cohort ($n=1{,}024$, 160 episodes, 63 runs).
Five closure rules were frozen before data generation; the v2 cohort
and its tables were committed before any v2 manuscript text was
written. One frozen-document miscount (``42 runs'' for a grid the
protocol itself defines as 63) was corrected ex ante in the engine
notes and is disclosed here.

\subsection{The frozen arms and closure rules}\label{sec:retire-protocol}
\paragraph{Retention levels.} (1)~\texttt{full-replay}: resubmit all
history. (2)~\texttt{masking}: identity-only context, every non-current
read fetched on demand---the strongest simple baseline.
(3)~\texttt{summary-reset}: periodic summary/reset every 4 episodes
with a window budget. (4)~\texttt{bounded-hier}: bounded pages.
(5)~\texttt{verified-board}: a present-state row per known gate.
(6)~\texttt{frontier}: verified-frontier retirement---retain only the
working set, boundary neighbors, dependency footprint, and pending
obligations, as a typed capsule (160\,B base, 16\,B per retained
binding actually held, 48\,B per pending obligation, 64\,B aggregate
invariants). (7)~\texttt{frontier-adaptive}: frontier plus a frozen
demand-driven eviction rule (bindings unused for 8 episodes and opened
$\ge$8 episodes ago are evicted outside the exempt base)---the
\emph{adaptive-coordination} candidate.
\paragraph{Closure rules (all frozen ex ante).}
\emph{Null closure}: frontier is promotable only if it cuts regime-O
cost per accepted run by $\ge$20\% versus masking pooled, with no
family stratum regressing by more than 10\%.
\emph{Attribution closure}: the scheduler factor is byte-neutral by
design; a byte leak exceeding 1\% of the grand mean voids all
scheduler results.
\emph{HACT closure}: grouped layouts are promotable only as a
conditional component---a $\ge$5\% win in the dense stratum plus a
stated density decision rule with reported regret.
\emph{Adaptive-coordination closure}: \texttt{frontier-adaptive} is
promotable only at equal acceptance, zero audit failure, $\ge$10\%
regime-O reduction versus frontier, and no family regression above 5\%.
\emph{Audit closure}: any arm with A3 defects in its audit tail is
reported as audit-failing and cannot be promoted.

\subsection{Results}\label{sec:retire-results}

\begin{table}[t]\centering\small
\caption{Frozen v2 cohort at $S=\mathrm{fixed}$: 84 cells $\times$ 7
retention policies. Regime-B bytes and regime-O cost are per accepted
run; ops are on-demand retrieval operations per accepted run; audit
defects (A3), staleness defects, and obligation losses are totals over
the 84 runs of each policy. Acceptance is the frozen contract: any
staleness defect or obligation loss rejects the run.}\label{tab:retire2-main}
\begin{tabularx}{\linewidth}{@{}l r r r r r r r@{}}\toprule
Policy & Accepted & Regime-B B/run & Regime-O B/run & Ops/run & A3 & Stale & Lost\\\midrule
\RetwofullreplayRow
\RetwomaskingRow
\RetwosummaryresetRow
\RetwoboundedhierRow
\RetwoverifiedboardRow
\RetwofrontierRow
\RetwofrontieradaptiveRow
\bottomrule\end{tabularx}\end{table}

\paragraph{The conversion rule is necessary: lossy conversion fails,
typed conversion does not.} Periodic summary/reset---the cheapest
policy on submitted bytes---accepts only \RetwoSummaryResetAccepted{}
runs: 109{,}761 stale reads and 385 lost obligations at the new scale,
the same structural failure v1 measured (30{,}820/280 on its smaller
grid). The mechanism is exact: a summary reverts values to a belief
older than the working window, and obligations crossing a reset
boundary survive only through a lossy summary. Every policy that
converts before retiring---frontier, and at higher cost the verified
board and full replay---accepts 84/84 with zero staleness and zero
loss. Bounded pages accept 84/84 but pay $1{,}211$ ops/run in the
steady phase; masking accepts 84/84 and pays $1{,}875$.

\paragraph{The null closure passes in the round-trip regime, not the
byte regime---exactly as pre-registered.} Frontier retirement reduces
regime-O cost per accepted run by \RetwoNullReduction{} versus masking
(gate $\ge$\RetwoNullThreshold{}), with \emph{no} family stratum
regressing (worst stratum: \RetwoWorstFamily). In the bandwidth regime
the same comparison is only \RetwoBytesDelta{}---we repeat v1's honest
negative and claim no byte-cost win. The two regimes separate because
frontier's advantage is \emph{operations}: \RetwoOpsReduction{} fewer
retrieval round trips per accepted run pooled
(66.5\% on clustered, 43.0\% on
independent; global coupling leaves no gap because every read is a
current fact for both policies). Sensitivity across the frozen
round-trip constant: \RetwoSensK{} at 1{,}024\,B, \RetwoSensM{} at
4{,}096\,B, \RetwoSensXxl{} at 16{,}384\,B---the advantage grows with
the round-trip price and is positive across the entire pre-registered
sensitivity range.

\paragraph{The audit tail is clean for the promoted mechanism.} Across
all 1{,}176 runs, the five-episode capsule audit records zero A3
defective service for every arm (audit closure:
\RetwoAuditClosure{}): retained bindings match stream generations,
pending sets match the stream's true pending sets in both directions,
and no required read was served from defective state. The audit also
caught a reference bug during development---the ground-truth pending
set originally deduplicated obligation openings by first rather than
latest opening, misclassifying correctly-held re-opened obligations as
phantoms---which is disclosed in the repository history and fixed
before the frozen cohort was regenerated.

\paragraph{Adaptive coordination was evaluated and rejected.} The
demand-driven eviction variant (\texttt{frontier-adaptive}) matches
frontier's acceptance (84/84) and audit cleanliness, but its regime-O
reduction versus plain frontier is $\RetwoAdaptiveReduction{}$
(promotion gate $\ge$10\%): the frozen use-clock rule evicts almost
nothing the base frontier does not already evict, because the base
exemptions (working set, obligations, footprint) already cover nearly
everything the use clock would keep. Per the frozen closure it is
\RetwoAdaptivePromotable{}, and the manuscript's mechanism carries no
adaptive component. This is the paper's second honest negative after
v1's byte-regime null result, and it was decided by a pre-registered
gate, not by narrative preference.

\paragraph{The scheduler factor is cleanly separated.} The full
factorial confirms byte-neutrality by design: the scheduler main effect
on regime-B bytes is \RetwoSEffectOnBytes{} of the grand mean (leak
threshold 1\%; attribution closure \RetwoAttributionClosure{}), while
the makespan interaction is large---adaptive batching cuts makespan
rounds by \RetwoMakespanReduction{} (24{,}471 to 3{,}070 rounds/run
pooled). Factor contrasts, not arm pairs, carry this claim: every
$R\times S$ combination was run in every cell.

\begin{table}[t]\centering\small
\caption{Family strata at $S=\mathrm{fixed}$: regime-B bytes and
regime-O cost per accepted run, and retrieval operations per accepted
run, masking versus frontier. Global coupling leaves no operations gap
(every read is a current fact); the operations advantage concentrates
where reads cross the frontier.}\label{tab:retire2-stratum}
\begin{tabularx}{\linewidth}{@{}l r r r r r@{}}\toprule
Family & Masking B/run & Frontier B/run & Regime-O cut & Masking ops & Frontier ops\\\midrule
\RetwoStratumClusteredRow
\RetwoStratumIndependentRow
\RetwoStratumGlobalRow
\bottomrule\end{tabularx}\end{table}

\paragraph{The operations advantage persists at four times the
horizon.} In the long-horizon sub-cohort ($n=1{,}024$, 160 episodes),
frontier cuts retrieval operations by \RetwoLongOpsReduction{} and
regime-O cost by \RetwoLongCostOReduction{} versus masking at equal
acceptance (63/63), while regime-B bytes remain within a few percent.
The advantage does not decay with horizon; if anything the clustered
stratum widens (66.5\%$\to$61.9\% operations cut at the pooled gate,
but 42.5\% regime-O reduction versus 39.8\% at the short horizon).

\begin{table}[t]\centering\small
\caption{Long-horizon sub-cohort ($n=1{,}024$, 160 episodes, seven
seeds, $S=\mathrm{fixed}$). Ops and costs are per accepted run. The
operations advantage persists at $4\times$ horizon; bytes stay
comparable.}\label{tab:retire2-long}
\begin{tabularx}{\linewidth}{@{}l l r r r r@{}}\toprule
Family & Policy & Accepted & Ops/run & Regime-O B/run & Regime-B B/run\\\midrule
\RetwoLongClusteredmaskingRow
\RetwoLongClusteredfrontierRow
\RetwoLongClusteredfrontieradaptiveRow
\RetwoLongIndependentmaskingRow
\RetwoLongIndependentfrontierRow
\RetwoLongIndependentfrontieradaptiveRow
\RetwoLongGlobalmaskingRow
\RetwoLongGlobalfrontierRow
\RetwoLongGlobalfrontieradaptiveRow
\bottomrule\end{tabularx}\end{table}

\paragraph{HACT remains a conditional component, now with a decision
rule that costs nothing.} The evidence-layout ablation (held-out half,
order learned on the training half, flat steel-man at identical record
cost) reproduces v1's shape at the new scale: the learned tree wins
the dense global stratum (169{,}408 vs 1{,}936{,}704 held-out B/cell,
an 87.1\% cut) and loses the sparse strata. The frozen density
rule---select the tree iff measured training-half update density
$\ge$0.5---matches the per-cell oracle on every one of the 84 cells
(\RetwoDensityRegret{} regret). HACT therefore stays what v1's closure
made it: a conditional serialization component with an observable
selection signal, not the paper's center. Section~\ref{sec:substrate}
retains the component development, condensed, with its scope
labels unchanged; the full derivations now live in the appendices.

\begin{table}[t]\centering\small
\caption{Evidence-layout ablation by family (pooled over sizes and
seeds; held-out half). The density rule picks the tree exactly when it
wins.}\label{tab:retire2-evidence}
\begin{tabularx}{\linewidth}{@{}l r r r@{}}\toprule
Family & Flat B/cell & HACT B/cell & Tree delta\\\midrule
\RetwoEvClusteredRow
\RetwoEvIndependentRow
\RetwoEvGlobalRow
\bottomrule\end{tabularx}\end{table}


\subsection{Pre-specified exploratory analyses (post-data, labeled as such)}\label{sec:retire-exploratory}
The five closures are the confirmatory results. The analyses in this
subsection were specified and committed \emph{after} the cohort was
frozen but \emph{before} they were run
(\texttt{docs/\allowbreak protocols/\allowbreak EXPLORATORY\_\allowbreak ANALYSES\_\allowbreak v2.md}); they are
exploratory and every number below is labeled as such. One amendment
(E6, the bounded-hier comparator) was itself prompted by observing
E1's pre-specified output; its motivation is disclosed in the protocol
document and summarized here.

\paragraph{E1/E5: the sensitivity sweep becomes a decision rule.}
Proposition~\ref{prop:breakeven} (Section~\ref{sec:sufficiency})
replaces the three-point sensitivity with a closed-form break-even:
with per-accepted means $(B_f,B_m,o_f,o_m)$, frontier satisfies the
20\% gate at round-trip price $c$ iff
$c\ge c_{0.8}=(0.8B_m-B_f)/(o_f-0.8o_m)$. On the frozen cohort
$B_f<B_m$ in every stratum, so $c_{0.8}=0$ pooled and per family:
frontier's byte advantage alone meets a 15\% gate before any
round-trip pricing, and every unit of price widens it
(pooled ratio $0.920\to0.588$ as $c$ goes $0\to32{,}768$\,B). The
gate is \emph{already satisfied at $c=0$}, which sharpens---and
honestifies---the earlier ``no byte win claimed'' phrasing: frontier
beats masking on bytes too, by 8.0\% pooled, just below the frozen
20\% byte threshold that the v1 protocol used. The 20\% byte gate
never fired in v2 because v2's null closure was pre-registered on
regime O, not regime B.

\paragraph{E2: distributional robustness.} Paired over the 84 cells
at $S=\mathrm{fixed}$: frontier's regime-O cost is below masking's in
\ExpFrontierMaskingSignO{} and its regime-B bytes below masking's in
\ExpFrontierMaskingSignB{} (exact two-sided sign tests). Per-seed
ops reduction ranges \ExpSeedOpsMin{}--\ExpSeedOpsMax{} and per-seed
regime-O reduction \ExpSeedOMin{}--\ExpSeedOMax{} across the seven
seeds: no seed reverses either direction.

\paragraph{E6: the bounded-recency comparator, analyzed head-on.}
E1's table exposed that \texttt{bounded-hier}---a 64-gate recency
window---is cheaper than frontier at the frozen 4{,}096\,B price in
three of four strata (bytes $0.168\times$ pooled), and the frozen
closure never compared them. The amendment was therefore written
before analysis. E6c: frontier beats bounded-hier in
\ExpFrontierBhSignO{} on regime-O and loses on bytes
(\ExpFrontierBhSignB{}); masking is dominated by bounded-hier in
\ExpMaskingBhSignO. E6a: the break-even price at which frontier
becomes cheaper than bounded-hier is \ExpBhBreakEvenPooled{} pooled
(\ExpBhBreakEvenIndependent{} independent; \ExpBhBreakEvenGlobal{}
global---bounded-hier wins at every price there). E6d (exploratory
runs): at the long horizon the picture splits by family
(Table~\ref{tab:retire2-long}). E6b is the mechanism-level
differentiator: bounded-hier's window leaves \ExpBhDueExposed{} of
due obligations and \ExpBhFootprintExposed{} of footprint reads
outside its retention, versus frontier's structural 0\%---its
correctness and cost both ride on the retrieval layer being free and
reliable, which the engine grants and a deployment may not. The
deployment recommendation
(Section~\ref{sec:deploy}) is updated accordingly: at prices below
the break-even, a deployment with reliable retrieval should prefer
the bounded window; the conversion rule remains the boundary of what
is safe without it.

\paragraph{E3: a real trace, replayed.} This repository's own git
history (\ExpTraceCommits{} non-merge commits; registry of
\ExpTraceRegistry{} file paths) is replayed through the same engine.
E3a (evidence layouts, real change sets only): the trace's
train-half density is \ExpTraceDensity{}, the density rule selects
the flat ledger, and the oracle agrees---\ExpTraceRuleRegret{} (the
tree would have cost $1.71\times$). The rule's first nonzero-density
real test passes. E3b (retirement, hybrid: real invalidation sets
with the frozen synthetic obligation/footprint overlay, disclosed as
such): frontier cuts operations \ExpTraceOpsCut{}
(\ExpTraceOpsRange{} across seeds) and regime-O \ExpTraceOCut{} with
\ExpTraceAcceptance{}---above the 20\% operations gate pooled, with
the pooled regime-O cut at 19.6\% marginally under it. Reported
exactly as measured: the real trace's burstier change sets compress
frontier's advantage relative to the frozen families, and the
live cohort remains the only arbiter of whether that survives real
retrieval.

\begin{table}[t]\centering\small
\caption{Pre-specified exploratory: bounded-hier at the long horizon
($n=1{,}024$, 160 episodes, seven seeds, $S=\mathrm{fixed}$; E6d),
beside the frozen frontier rows of Table~\ref{tab:retire2-long}.
Bounded-hier wins independent and global; frontier wins clustered.
Exploratory, not confirmatory.}\label{tab:bh-long}
\begin{tabularx}{\linewidth}{@{}l l r r r@{}}\toprule
Family & Policy & Ops/run & Regime-O B/run & Regime-B B/run\\\midrule
clustered & bounded-hier & \ExpLhBhClusteredOps & \ExpLhBhClusteredO & \ExpLhBhClusteredB\\
independent & bounded-hier & \ExpLhBhIndependentOps & \ExpLhBhIndependentO & \ExpLhBhIndependentB\\
global & bounded-hier & \ExpLhBhGlobalOps & \ExpLhBhGlobalO & \ExpLhBhGlobalB\\
\bottomrule\end{tabularx}\end{table}

The exploratory results do not change the paper's confirmatory
claims; they change its deployment advice and its honesty about the
bounded-recency alternative. The live cohort's arms
(Section~\ref{sec:live-protocol}) are unchanged: masking,
summary/reset, and frontier remain the arms the frozen closures
separate, and bounded-hier enters the live protocol only as a fourth
arm if its retrieval-exposure profile can be instrumented there.

\subsection{What the mechanism study does and does not
establish}\label{sec:retire-limits}
The conversion rule and the regime-conditional payoff are structural
predictions about policies, not measurements of agents. The audit
checks a policy's capsule against stream ground truth in a
deterministic engine; a live agent's summaries, retrievals, and
obligations are not instrumented by it. The $4{,}096$\,B
round-trip constant is a frozen accounting choice, not a measured
latency; the sensitivity range brackets it but does not replace a live
measurement. Byte numbers are the frozen accounting of a frozen model
of information flow; no hosted model was called, no token or billing
claim is made, and the BPE and physical-WAN gaps of
Section~\ref{sec:eval} remain open. The live cohort of
Section~\ref{sec:live-protocol} pre-registers operations as its
primary cost measure precisely because that is where the mechanism
result predicts the payoff.

\section{Sufficiency of the Conversion Rule Within the Model}\label{sec:sufficiency}

The frozen cohort shows the conversion rule is \emph{necessary} at
mechanism level: every measured policy that retires before
converting fails the acceptance contract. This section proves the
complementary direction---\emph{sufficiency}---for the frontier policy
under the engine's semantics, so the paper's rule is
necessity-plus-sufficiency within the model rather than an empirical
regularity. The proposition is scoped to the deterministic engine and
to streams with reliable retrieval (every fetch succeeds and returns a
current value); the unreliable-retrieval boundary is exactly where the
rule's guarantees stop, and Section~\ref{sec:retire-results} measures
what bounded-hier's exposure looks like there.

\begin{proposition}[Frontier sufficiency]\label{prop:sufficiency}
Fix a finite episode stream $E_0,\dots,E_{T-1}$ satisfying the frozen
stream invariants: (i) every fact's payload is observed before or at
its first required read or at its due obligation's resolution episode;
(ii) every obligation opened at episode $o$ has its gate in the
required reads of episode $\min(o+W, T-1)$, where $W$ is the
obligation window; (iii) every required read of episode $e$ is either
a current fact of $e$ or fetchable. Then the frontier policy incurs
no staleness defects, no obligation losses, and resolves every
obligation opened at $o$ with $o+W\le T-1$.
\end{proposition}

\begin{proof}
Staleness. The frontier policy never stores a believed value: a gate
enters \texttt{known} only through a fresh fact or a successful fetch,
both of which set its generation to the current episode, and eviction
only removes entries. A required read is answered from the current
facts, from a retained entry, or by a fetch (invariant (iii)); the
first two are current by construction---retained entries are updated
by every observation of their gate---and the third is current by the
reliable-retrieval assumption. Hence no read is served stale.

Obligations. An obligation opened at $o$ is inserted into the pending
set at $o$ and is exempt from eviction (the frontier retains all
pending gates). By invariant (ii) its gate is a required read at
$r=\min(o+W,T-1)$. At $r$, the gate is a current fact, or is retained
(the pending exemption keeps it in \texttt{known} from opening), or is
fetched; in all cases the resolution condition of the engine (the gate
was read validly at its due episode) is met, so the obligation
resolves. No pending obligation is ever dropped by the frontier
policy, so no loss is recorded. For $o+W\le T-1$, $r=o+W\le T-1$, so
the resolution happens within the horizon.
\end{proof}

\begin{remark}
The two directions use disjoint evidence. Necessity is measured: every
lossy-conversion policy fails on the frozen grid
(Section~\ref{sec:retire-results}). Sufficiency is proved under
invariants the frozen stream generator guarantees by construction,
which is also why the 84/84 acceptance of frontier runs is a
\emph{check} of the implementation against the proposition rather than
independent evidence for it. What the model does not prove is the
reliable-retrieval assumption itself: in a deployment where re-fetches
can fail or return stale state, the conversion rule's safety reduces
to the retrieval layer's guarantees, and bounded-hier's measured
exposure (69.0\% of due obligations and 35.9\% of footprint reads
outside its window) quantifies how much a bounded-recency design
delegates to that layer.
\end{remark}

\paragraph{Retrieval-failure sensitivity.}
The reliable-retrieval boundary admits a closed-form sensitivity
analysis. The following proposition is a declarative computation over
the already-measured operation counts, not a new experiment: the
engine's fetches never fail, so no data exists at $p>0$ and none is
simulated here. It quantifies exactly where Proposition
\ref{prop:sufficiency}'s guarantee stops, and why the conversion rule,
not a looser policy, is what survives there.

\begin{proposition}[Degradation under unreliable retrieval]\label{prop:pfail}
Fix the stream invariants of Proposition~\ref{prop:sufficiency},
replacing invariant (iii) by: each fetch succeeds independently with
probability $1-p$ for $p\in[0,1)$, and a failed fetch leaves its
required read unvalidated---recorded as a staleness defect, never
zero-imputed. Let $R(\mathcal P)$ denote the number of fetch-dependent
required reads of policy $\mathcal P$ on the stream. Then:
\begin{enumerate}[leftmargin=*,itemsep=0pt]
\item[(i)] \texttt{full-replay} has $R=0$ and incurs no
retrieval-induced defects at any $p$;
\item[(ii)] $\mathcal P$'s expected retrieval-induced defects equal
$p\,R(\mathcal P)$, and its per-run defect-free probability is
$(1-p)^{R(\mathcal P)}$;
\item[(iii)] if $R(\text{frontier})<R(\text{masking})$, frontier's
expected defects are strictly below masking's at every $p>0$, and the
ratio is $p$-independent: the measured operations gap carries over
unchanged as a defect-rate gap;
\item[(iv)] \texttt{summary-reset}'s defect count is strictly positive
at $p=0$ (109{,}761 stale reads and 385 lost obligations occurred with
retrieval fully reliable), so no retrieval layer, however perfect,
makes it safe.
\end{enumerate}
\end{proposition}
\begin{proof}
(i) \texttt{full-replay} resubmits all history and issues no fetch.
(ii) Each fetch-dependent read fails independently with probability
$p$; the count of failed reads is binomial with mean $p\,R(\mathcal P)$,
and a run is defect-free exactly when all $R(\mathcal P)$ fetches
succeed. (iii) Immediate from (ii) by linearity:
$p\,R(\text{frontier})<p\,R(\text{masking})$ for $p>0$, and the ratio
$R(\text{frontier})/R(\text{masking})$ does not depend on $p$.
(iv) These defects arise from cross-boundary information drops inside
the summary itself, not from fetches. \qed
\end{proof}

Three quantitative consequences use the long-horizon operation counts
of Table~\ref{tab:retire2-long}. First, the degradation point against
full-replay: frontier's expected excess defects reach one per accepted
run at $p\approx1/284\approx3.5\times10^{-3}$ in the clustered family
and $p\approx1/8{,}943\approx1.1\times10^{-4}$ in the independent
family; the global family has $R=0$ and is structurally immune at
every $p$. More generally, a deployment tolerating at most $\epsilon$
expected defects per run requires $p<\epsilon/R(\text{frontier})$.
Second, the ordering against masking is preserved: frontier's pooled
fetch volume is \RetwoOpsReduction{} smaller, so its expected
retrieval-induced defects are $56\%$ of masking's at every failure
rate---the conversion rule converts an operations advantage into a
correctness advantage precisely where retrieval is unreliable, and the
break-even price of Proposition~\ref{prop:breakeven} is not needed for
this direction. Third, bounded-hier enters the same formula with
$R$ equal to its out-of-window read volume, which the cohort measured
as exposure fractions (\ExpBhDueExposed{} of due obligations,
\ExpBhFootprintExposed{} of footprint reads) rather than operation
counts; its safety delegates to the retrieval layer in exactly the
sense of item (ii), with an exposure larger than frontier's
structural zero for pending obligations.

The proposition also marks the boundary of what is claimed. A real
retrieval layer fails with correlation, staleness, and byzantine
responses, not independent Bernoulli trials; retries turn $p$ into
$p^{r}$ only under independent attempts. Neither refinement is
measured here. What the model does establish is structural: the
conversion rule makes the failure surface \emph{countable}---$R$ is a
measurable property of the typed capsule---where a summary's failure
surface is silent, and item (iv) shows the difference is not a matter
of retrieval quality at all.

\begin{proposition}[Monotone break-even rule]\label{prop:breakeven}
Let $B_f,B_m>0$ and $o_f,o_m\ge 0$ be per-accepted-run regime-B means
and per-accepted-run operation means of two policies, with
$o_f<o_m$. Define $R(c)=(B_f+c\,o_f)/(B_m+c\,o_m)$ for $c\ge 0$. Then
$R$ is monotone on $c\ge0$: decreasing iff $B_f/B_m>o_f/o_m$,
increasing iff $B_f/B_m<o_f/o_m$, constant iff equal; its endpoints
are $R(0)=B_f/B_m$ and $\lim_{c\to\infty}R(c)=o_f/o_m$. For any
$\rho\in(o_f/o_m,\;B_f/B_m)$ (nonempty exactly when $R$ is
decreasing),
\begin{equation}
R(c)\le\rho\iff c\ge c_\rho,\qquad
c_\rho=\frac{\rho B_m-B_f}{o_f-\rho\,o_m}>0.
\end{equation}
If $\rho\ge B_f/B_m$ the gate holds at $c=0$; if $\rho\le o_f/o_m$ it
holds at no $c$.
\end{proposition}

\begin{proof}
$R'(c)$ has the sign of $o_f B_m-o_m B_f=B_mB_m\,(o_f/o_m-B_f/B_m)$,
which is constant in $c$, giving monotonicity with the stated
directions and endpoints. Suppose $o_f/o_m<\rho<B_f/B_m$. Solving
$B_f+c\,o_f\le\rho(B_m+c\,o_m)$ for $c$ gives
$c\,(o_f-\rho o_m)\le\rho B_m-B_f$. Here
$\rho B_m-B_f>0$ since $\rho>B_f/B_m$, and $o_f-\rho o_m<0$ since
$\rho>o_f/o_m$; dividing by the negative coefficient reverses the
inequality, so the gate is equivalent to $c\ge c_\rho$ with
$c_\rho=(\rho B_m-B_f)/(o_f-\rho o_m)>0$. If $\rho\ge B_f/B_m=R(0)$,
monotone decrease gives $R(c)\le R(0)\le\rho$ at every $c\ge0$; if
$\rho\le o_f/o_m=\lim_{c\to\infty}R(c)$, then $R(c)>\rho$ at every
finite $c$. The case $o_m=0$ forces $o_f=0$, so $R$ is constant at
$B_f/B_m$ and only the first endpoint rule applies.
\end{proof}

\begin{remark}[Deployment use]
Proposition~\ref{prop:breakeven} turns the frozen sensitivity sweep
into a decision rule. A deployment pilot that estimates the four means
$(B_f,B_m,o_f,o_m)$ and the deployment's round-trip price $c$ obtains
an exact go/no-go for the frozen 20\% promotion gate: adopt frontier
iff $c\ge c_{0.8}$, where
$c_{0.8}=(0.8B_m-B_f)/(o_f-0.8\,o_m)$ whenever
$o_f/o_m<0.8<B_f/B_m$. On the frozen cohort the pilot must actually
price its round trips: frontier's byte advantage alone ($8.0\%$
pooled) does not meet a 15\% or 20\% gate, so
$c_{0.8}\approx1{,}364$\,B pooled, $c_{0.85}\approx658$\,B pooled,
and $c_{0.8}$ is never met in the global family, where both policies
are fetch-free and the ratio is constant at $0.924$
(Section~\ref{sec:retire-exploratory}).
\end{remark}

\section{Verification Substrate and Serialization}\label{sec:substrate}
The retirement study of Section~\ref{sec:retire} asks when history can
be discarded; this condensed section defines the verification
substrate that makes the discarded history reconstructable and
auditable---the stable identities, generations, and evidence layouts a
typed capsule refers to. The conversion heuristic of
Section~\ref{sec:sufficiency} consumes this substrate; it does not
depend on any particular layout optimizer. These component results
carry their v5.1 scope labels unchanged; they are no longer the
paper's center, and the new ablation (Table~\ref{tab:retire2-evidence})
now bounds when their layout choice matters at all. Full derivations
that formerly occupied several sections here are consolidated; the
exact-DP formulation and its proof now live in
Appendix~\ref{app:exactdp}, and the block-backbone ablation remains in
Appendix~\ref{app:block}.

\subsection{Model and trust boundary}\label{sec:model}
\paragraph{Semantic registry, execution, and display.}\label{sec:registry}
Let $\R=(g_0,\ldots,g_{n-1})$ be a nonempty approved registry of unique check
IDs. The candidate identity is
\begin{equation}e=( \text{epoch},\ \text{snapshot},\ \text{checker},\ \text{environment},\ \text{registry}).\end{equation}
The snapshot covers source and protected acceptance inputs. The checker and
environment are explicitly bound; an epoch increases even after an ABA return to
identical source bytes. Only a trusted runner submits final \PASS{}, \FAIL{}, or
\UNKNOWN{} records. Missing, stale, skipped, incompatible, or timed-out outcomes
are not passes.

The model splits its state into two disjoint layers.
\emph{Immutable check identities} are the registry $\R$, the candidate
identity $e$, and the acceptance predicate each check declares; they are fixed
at publication and are never rewritten by execution, layout choice, or
learning. \emph{Mutable execution state} is everything a run may change: the
result tokens, their generation counters, epoch-local timestamps, and the
display positions. An execution permutation $\sigma_e$ determines when checks
run. A layout $L=(\pi,T)$ independently maps tree positions to semantic
registry indices. For a node covering positional interval $[i,j]$, its semantic
coverage is
\begin{equation}G_L(i,j)=\{g_{\pi(i)},\ldots,g_{\pi(j)}\}.\label{eq:coverage}\end{equation}
The stored layout must contain an exact permutation, and every internal node's
children must partition its positional interval without gaps or overlaps. The
task graph is outside this mathematical tree; one worker can discharge many
checks, and one integration check can cover several workers.

The acceptance predicate is independent of both permutations:
\begin{equation}A(e)=\mathbf 1\{\text{every }g_i\in\R\text{ has trusted, matching, current \PASS{} evidence for }e\}.\label{eq:accept}\end{equation}
An observed fail permits rejection before all checks finish. It does not assert
coverage of the unexecuted remainder. Acceptance is relative to the approved
predicates, not proof of all program behavior.

\paragraph{Packet and lifecycle costs.}\label{sec:packet}
Each internal node of arity $d$ emits a fixed-width packet
\begin{equation}c(d)=H+Sd,\qquad 2\le d\le b=\left\lfloor (C-H)/S\right\rfloor.\label{eq:price}\end{equation}
Here $H=512$, $S=320$, and $C=3072$ bytes, so $b=8$. A mandatory field that
exceeds its slot is rejected, not silently truncated. The cap applies to one
certificate packet, not to all project memory, an entire registry manifest, or
every LLM prompt. A manifest is a referenced control-plane artifact; its measured
size is charged on migration.

Let $W_{e,1},\ldots,W_{e,q_e}$ be the semantic IDs whose final records become
available between successive refreshes. For a strict new candidate epoch, the
observed cost is
\begin{equation}\ell_e(L)=W(L)+\sum_{a=1}^{q_e}\sum_{v\in\mathrm{int}(T)}c(d_v)\,\mathbf 1\{W_{e,a}\cap G_L(v)\ne\varnothing\},
\qquad W(L)=\sum_{v}c(d_v).\label{eq:lcost}\end{equation}
The first term constructs the full initially unknown tree. A successful strict
epoch runs the whole registry. We separately study persistent-state refresh
traces that omit $W(L)$ at each tick, as in a cache with correctly specified
incremental dependencies; each topology installation is still cold and
explicitly priced. These traces evaluate the controller's cost model; they are
not the strict candidate-gate timing experiment, and they do not demonstrate
safe dependency inference from arbitrary code.

\paragraph{Why more than invalidation or pairwise similarity is needed.}\label{sec:why}
A classical result shows that marginal check activation cannot identify the best
aggregation tree under joint invalidation (Appendix~\ref{app:sep}), and even the
full initial invalidation distribution cannot identify optimal progressive
verification. One further construction bounds what pairwise statistics can see:

\begin{proposition}[Pairwise affinity is insufficient for full-wave scoring]\label{prop:pairwise}
There exist two distributions with identical singleton and pairwise co-activity
probabilities, hence identical Jaccard affinities, but different activation
probabilities for the same three-check subtree.
\end{proposition}
\begin{proof}
On IDs $\{0,1,2\}$, let $P$ be uniform on $\varnothing,\{0,1\},\{0,2\},\{1,2\}$
and $Q$ uniform on $\{0\},\{1\},\{2\},\{0,1,2\}$. Every singleton has
probability $1/2$, every pair has joint probability $1/4$, and its Jaccard
affinity is $1/3$. Yet the probability that the triple is hit is $3/4$ under $P$
and $1$ under $Q$. Thus a packet emitted on that interval has different expected
cost. \qed
\end{proof}
The implication for the compiler is specific: a pairwise clustering
objective cannot substitute for scoring recorded completion hyperedges.

\subsection{Learning layouts without changing the contract}\label{sec:compiler}
\paragraph{Practical catalogue and full-wave selection.}\label{sec:catalogue}
Training waves define counts $N_i$, $N_{ij}$ and affinity
\begin{equation}a_{ij}=\frac{N_{ij}}{N_i+N_j-N_{ij}},\label{eq:affinity}\end{equation}
with zero for an unseen pair. Affinity orders semantically related checks
canonically. All semantic IDs, including disconnected and unseen checks, remain
included exactly once.

Given the learned order, group consecutive leaves into blocks of at most eight,
promote a singleton unchanged, and repeat upward. This bottom-up compact balanced
construction does not call a tree optimizer. The catalogue retains the incumbent
first, includes original-order fixed8 when distinct, and proposes learned-order
fixed8. Identical layouts are deduplicated. A separate validation cohort scores
complete completion waves, including initial construction; ties retain the
incumbent. The API returns an advisory choice, not deployment authorization. The
registry and packet cap cannot change during selection.

For any proposal $\pi$, training episodes supply interval weights
\begin{equation}r_\pi(i,j)=1+\frac{1}{m}\sum_{e=1}^{m}\sum_{a}\mathbf 1\{W_{e,a}\cap G_L(i,j)\ne\varnothing\}.\label{eq:weights}\end{equation}
Pairwise affinity proposes an order; it does not replace the full-wave cost. The
practical catalogue contains at most three candidates and uses no DP. Tree
construction is $O(n)$ after ordering. The reference implementation retains
spectral ordering and exact fitting for research comparisons; neither is invoked
by the practical default. The exact conditional tree synthesis---its dynamic
programming recurrence, exactness theorem, and proof---is retained unchanged in
Appendix~\ref{app:exactdp} for reference; the practical default never calls it.

\begin{proposition}[Validation after learned ordering]\label{prop:valid}
Condition on a training set that fixes $K$ complete layouts. For $m$
independent identically distributed validation epochs with each layout cost in
$[0,B]$, let $\widehat L$ minimize empirical validation cost. With probability at
least $1-\delta$,
\begin{equation}J(\widehat L)\le\min_{L\in\mathcal C}J(L)+2B\sqrt{\frac{\log(2K/\delta)}{2m}}.\label{eq:valid}\end{equation}
\end{proposition}
\begin{proof}
Conditionally, the catalogue is fixed. Hoeffding's inequality and a union bound
over $K$ layouts bound every empirical mean error by the displayed radius
without its factor two. Inserting empirical costs on either side of the empirical
minimizer yields the result. \qed
\end{proof}
The guarantee is only relative to the fixed catalogue. It neither justifies
choosing among learned permutations on their training errors alone nor covers
distribution shift. We do not report this potentially loose bound as an empirical
confidence interval; reported intervals use independent experimental seeds.

\paragraph{Compiler identity and scope.}\label{sec:compiler-scope}
The practical default has no hidden microtree optimization. Exact DP is a
reference and optional offline candidate (Appendix~\ref{app:exactdp}). The
block-backbone proposition and implementation remain in Appendix~\ref{app:block}
solely as an ablation; its restricted-class exactness does not override the
observed simpler-baseline advantage. At 1{,}024 checks, order learning still
takes seconds in our measurement, so layouts should be cached and compiled away
from the agent's per-action path. Validation of padded costs does not optimize
compressed archives or model tokenization; adoption must measure the actual
delivery representation.

\begin{algorithm}[t]\caption{Deployable evidence control; inference and checking remain separate}\label{alg:control}
\begin{algorithmic}[1]
\For{each build request}
\State Build context only from the semantic registry
\State Check out the frozen, identified source snapshot
\State If required, select layout offline via the frozen rule
\If{layout changes}
\State Revoke tickets, increment generation, remap stable IDs, rebuild and price the installation
\EndIf
\State Receive trusted completion records for the fixed check schedule
\State Reject stale, conflicting or malformed records; aggregate bounded packets
\State Publish only the complete current contract result; otherwise retain \FAIL/\UNKNOWN
\State Expose a root-status card; retrieve diagnostic packets only when needed
\State Score the observed event and append its cost vector to the bounded window
\EndFor
\end{algorithmic}\end{algorithm}

\subsection{Practical online selection: price the window, not the promise}\label{sec:online}
Freeze a compiled catalogue $\mathcal C=\{L_1,\ldots,L_K\}$. At event $t$, choose
a layout before observing the current completion wave; afterwards score that same
wave under every catalogue member. This is full-information feedback only because
the check schedule is held fixed. It is not counterfactual feedback for arbitrary
agent behavior.

A different layout is installed by serializing its canonical manifest and
reconstructing its complete certificate tree:
\begin{equation}M_{ij}=\begin{cases}0,&L_i=L_j,\\ |\mathrm{manifest}(L_j)|+W(L_j),&L_i\ne L_j.\end{cases}
\qquad L_T=\sum_{t}\ell_t(A_t)+\sum_{t=2}^{T}M_{A_{t-1},A_t}.\label{eq:mig}\end{equation}
The initial installation is separate. These byte prices match the evaluated
uncompressed migration protocol. A deployment using compression must measure its
own priced wire representation; the DP objective is not automatically an optimum
for a compressed log.

\paragraph{The primary practical rule.}
Let $\bar\ell_{t,i}$ be the mean cost on the last $\min(w,t-1)$ observed events,
and $j_t=\argmin_i\bar\ell_{t,i}$ with deterministic ties. Starting from the
calibrated incumbent $a$, use the frozen rule
\begin{equation}A_t=\begin{cases} j_t,& w(\bar\ell_{t,a}-\bar\ell_{t,j_t})>M_{a,j_t},\\ a,&\text{otherwise},\end{cases}\qquad w=64.\label{eq:window}\end{equation}
With no history, keep $a$. A tie in estimated savings does not trigger migration.
During warm-up, the forecast horizon remains $w$ even if fewer events were
observed, matching the original baseline. This is an empirical break-even
forecast, not proof that past savings repay a future switch. It can lose under
abrupt reversal; preserving the frozen layout is always a legitimate alternative.

The runtime API enforces \texttt{choose()} before \texttt{observe()}, bounds
history to 64 vectors, rejects nonfinite costs, and checkpoints only at event
boundaries. Replaying all 80 archived streams, with periodic checkpoint/resume,
must reproduce the original baseline's full action and byte-cost paths exactly.
The exercise validates implementation equivalence, not a new performance dataset.

\paragraph{Why fixed-share is no longer primary.}
On the archived stationary, long-shift and rapid-shift streams, the coupled
fixed-share policy is more expensive than the priced window. Its standard
tracking analysis therefore moves to Appendix~\ref{app:fixedshare}, with all
negative results retained. We do not manufacture an adversarial workload merely
to make that policy win. Neither the appendix bound nor the empirical window
result establishes general superiority over other switching-cost policies.

\subsection{Stable evidence through permutation and migration}\label{sec:soundness}
A result token carries $(\text{epoch},g_i,b_i)$, where $b_i$ binds the semantic
ID to the registry, source snapshot, checker, and environment. It does not carry
a display position as its identity. A layout-specific kernel maps stable IDs
to current positions. During migration it rebuilds from canonical trusted records;
it never transplants an ordinal cache entry into a different permutation.

Publication tickets additionally carry a monotone layout generation and layout
hash. A layout switch revokes the issued ticket, advances the generation, and
forces a full recheck of the affected subtree. The registry, candidate identity,
and acceptance predicates never change through a layout-only
change, but a new candidate epoch fences it. Duplicate identical evidence is
idempotent. Conflicting final records in one epoch revoke authorization: a
recorded fail remains a fail; otherwise the obligation becomes unknown.
Resolution requires a new epoch, not a more persuasive agent summary.

\begin{theorem}[Layout-independent local publication soundness]\label{thm:soundness}
Assume complete approved predicates and dependency identities, a trusted checker
executing immutable identified inputs, collision-resistant binding, and atomic
local epoch/generation checks. For any sequence of valid layout permutations,
any advisory learning outputs, and any delivered evidence order, authorization
implies Equation~\eqref{eq:accept}. If each checker is sound for its declared
predicate, those predicates hold on the authorized candidate.
\end{theorem}
\begin{proof}
A canonical pass requires a current token for the same semantic ID and immutable
identities. Permutation is a bijection, so rebuilding moves each such record to
exactly one position without changing its meaning. Exact child partitioning
preserves coverage and pass/fail/unknown counts by induction up the tree. The
root can be authorized only with $(n,0,0)$ and matching current registry, epoch,
generation, and layout. Old tickets fail those checks. Conflicts prevent
authorization. No choice made by the learner writes trusted outcomes or changes
the registry. \qed
\end{proof}
This is a conditional local guarantee, not remote attestation or distributed
consensus. An incomplete suite, compromised checker, omitted dependency, or
nondeterministic external service can invalidate a broader correctness claim. The
CLI's file copy is not an OS sandbox. An external merge/deployment service must
atomically compare the verified candidate and integration-head identities at
commit; a report file alone is not that transaction.

\subsection{Implementation and cost accounting}\label{sec:impl}
The artifact adds \texttt{practical.py}, \texttt{deployment.py},
\texttt{token\_budget.py} and \texttt{tls\_probe.py}. The practical CLI produces
a balanced catalogue without calling either exact or block DP. The historical
exact and block compilers remain separately callable. The existing stable-ID
guard, priced window, codec and Future Code bridge are retained. The strict CLI
compares protected input sets and hashes, detecting added tests or configuration
files as well as changed and deleted ones. Schema-2 contracts bind the approved
checker and environment and refuse silent reuse of legacy contracts. Source
manifests are checked before and after snapshotting and verification; mutation
of an observed snapshot prevents a pass.

The bridge runs as an ordinary subprocess acceptance gate in the unmodified
Future Code v1.2 runtime. The supervisor still owns worker leases, private
attempts, bounded contexts, and serialized integration. Non-pass exits cause
candidate rejection. Detailed evidence and binary certificate packets live
outside candidate write scopes; only a bounded JSON result card returns to the
worker. The plug-in accepts a precompiled layout. The online selector and
migration guard are separately integrated and stress-tested as a library path; a
live, self-retuning hosted deployment was not evaluated.

A different scalar objective could combine these quantities only with declared
unit conversions and observable counterfactual costs. This implementation does
not optimize unknown model latency or billing. The exact-byte claim also depends
on the stated padded format; compression defines a different, potentially
nonseparable objective for which this DP is not claimed exact.

Costs are accounted in separate ledgers so byte savings are never converted into
unsupported latency or token claims:

\begin{table}[t]\centering\small
\caption{Separate ledgers avoid converting byte savings into unsupported latency or token claims.}\label{tab:ledgers}
\begin{tabularx}{\linewidth}{@{}lXX@{}}\toprule
Quantity & Instrument & Ledger\\ \midrule
Serialized certificate bytes & Sum over packets with fixed per-node header and slot widths & Bytes; never tokens\\
Wall-clock durations & Monotonic clock around whole-run driver and per-stage timers & Seconds; never bytes\\
Tool calls & Counted supervisor-side & Counts\\
Check latencies & External driver timeouts are host-configured and enforced per tool call & Reported separately\\
Migration bytes & Canonical layout manifest plus forced full rebuild & Yes, online traces\\
Fit and scoring time & Monotonic-clock measurements & Reported separately\\
Preparation, copying, checking, aggregation & Wrapper timers and outer call duration & Reported separately\\
Active harness episode & Supervisor startup through completion and shutdown & Reported separately\\
Prompt bytes and construction time & Captured actual HTTP messages and context builder & Reported; not tokens\\
Hosted token usage and API billing & No provider inference was run & Unknown, never zero-imputed\\
\bottomrule\end{tabularx}\end{table}

\section{Service-Matched Delivery and Admission}\label{sec:delivery}
\subsection{Do not conflate three consumers}\label{sec:consumers}
The local aggregator receives bounded packets; a logging collector may receive
their complete stream; a supervisor normally receives a status projection and
retrieves details selectively. Define $\psi(e,c)$ as canonical status text
containing semantic identity, counts and local authorization, but no display
positions. For the same trusted semantic report,
\begin{equation}\psi_{L_1}(e,c)=\psi_{L_2}(e,c).\label{eq:rootid}\end{equation}
Its byte length and token count under any fixed deterministic tokenizer are
therefore identical. This identity is not a new theorem about tokenization. It
supplies a negative control: layout savings cannot reduce a layout-independent
root input. The card is a projection of a trusted report, not authority to accept
arbitrary worker JSON.

Diagnostics preserve full header/card records and page whole packets under a
16{,}384-byte limit. A packet that cannot fit is rejected, not truncated. For a
real supplied tokenizer $\tau$ and explicitly reserved framing budget $r$, the
new admission API additionally requires
\begin{equation}|p|_{\mathrm{UTF8}}\le C_{\text{bytes}},\qquad |\tau(\mathrm{decodeUTF8}(p))|+r\le C_{\text{tokens}}.\label{eq:admit}\end{equation}
The implementation retokenizes each full proposed page, including its prefix.
Token counts are not assumed additive across concatenations. Greedy packing
guarantees the tested limits, not a minimum number of pages. Toy counters in API
tests are not BPE measurements. The optional measurement command requires actual
tiktoken software and vocabulary~\cite{tiktoken}, fingerprints the encoding, and
records per-page content hashes. Missing encodings produce \UNKNOWN{} and null
counts; there is no byte-ratio fallback. Provider framing and billed usage
require a separately measured adapter, not a guessed constant.

\subsection{Encoding and the simplest sufficient service}\label{sec:encoding}
Canonical unpadded JSON, independently compressed packets, and whole-epoch
DEFLATE preserve their input records with different operational boundaries.
Whole-epoch compression requires the completed archive, whereas per-packet
compression permits incremental export. The padded objective is not an exact
optimizer of either compressed bytes or tokens. Practical selection is advisory
until its gain is measured with the deployed codec.

For a trusted consumer needing only final status, a root endpoint or flat ledger
can be much cheaper than the intermediate stream. This changes the diagnostic
service, not necessarily the underlying check set. Comparing that smaller service
with an exhaustive stream cannot demonstrate that a tree optimizer is efficient
or necessary. The strict checker, data representation, and delivery demand must
each be named explicitly.

\subsection{Adoption cost is not marginal layout cost}\label{sec:adoption}
Consider $q$ copies of the same workload and delivery service through one
nonoverlapped bottleneck. Let $B_0,B_1$ be codec-matched application bytes per
copy, $R>0$ effective application goodput, $E\ge 0$ additional
encoding/decoding work per copy, and $C\ge 0$ consistently amortized one-time
incremental work. Common setup cancels only when actually shared. The conditional
net benefit is
\begin{equation}d=(B_0-B_1)/R-E,\qquad \Delta(q)=qd-C.\label{eq:netbenefit}\end{equation}
A strictly profitable integer copy count exists only when $d>0$, and then
\begin{equation}q_{\min}=\lfloor C/d\rfloor+1.\label{eq:qmin}\end{equation}
This is accounting, not a new networking or regret theorem. A zero gain retains
the baseline. The implementation uses decimal arithmetic at equality and rejects
mismatched service IDs. Missing $C$, $E$ or $R$ returns \UNKNOWN{}, not zero. An
explicit zero is permissible only as a declared assumption or measurement.

When both variants already use the same protected gate, $C$ is extra layout work,
not the full strict-wrapper cost. When adopting strict verification from ordinary
testing, its additional cost must be justified independently; the services differ
and the planner returns \texttt{INCOMPARABLE}. Repeating exports amortizes one
candidate's preparation, not the preparation of future new candidates. Parallel
links, computation overlap, cached diagnostics, queueing and different connection
policies require direct critical-path measurement. High RTT alone does not make
byte savings larger if both exports pay the same setup and request count.

\subsection{Remote measurement and security boundaries}\label{sec:remote}
The historical transfer experiment is receiver-paced loopback, not a WAN. The new
opt-in measurement kit uses verified TLS server certificates and hostname
checking; upload and nonlocal bind require explicit switches. Frames and
decompressed archives are bounded, and acknowledgments check received length and
digest. Local tests exercise cold and reused TLS sessions, rejected hostnames,
untrusted roots, malformed acknowledgments and truncated records.

These tests validate a protocol path only. They provide no physical-WAN
throughput, jitter, packet-loss or retransmission measurement. The supplied
protocol records host and path provenance, codec, connection policy and matched
payloads, and requests separately observed RTT/goodput and transport diagnostics
where available, consistent with the distinctions in TCP throughput
testing~\cite{rfc6349}. Fields that the minimal client cannot observe remain
null. There is no receiver-side artificial pacing in a qualifying WAN trial. No
public upload, account creation or credential search was performed.

TLS authenticates the configured server, not the checker's truthfulness. Protected
reference data, immutable candidate identity and local monotone state remain
required for Theorem~\ref{thm:soundness}. There is no remote attestation, Byzantine
consensus, client authorization service, non-repudiation or OS sandbox. Replacing
both evidence and its trusted reference still defeats a digest-only comparison.
The receiver is a bounded measurement utility to use behind an operator-controlled
firewall, not a production ingestion service.

\section{Evaluation of the Practical Revision}\label{sec:eval}
\subsection{Provenance and frozen questions}\label{sec:provenance}
The v5 contract was saved before the new runs. It fixes three sizes (128, 512,
1{,}024), three completion families (clustered, independent, global), and three
seeds (55011--55013). Each of the 27 workloads has 128 training epochs, 128
validation epochs and 256 held-out epochs. Random draws are sequential disjoint
samples; deterministic global observations naturally coincide. The catalogue is
fit only on training data, selected on validation, and reported on held-out data.
Cold strict-epoch construction is included in every objective. The nine
size/family summaries use seed means and observed ranges, not episode-pooled
significance tests.

A separate study re-aggregates twelve archived source candidates under three
layouts, producing 36 layout replays. All use the same original-order trusted
check records, the same executed sequence and batches of eight.
Average-linkage orders were learned on the earlier training cohort and frozen
before these candidates. None of these replays is a fresh checker execution,
independent code repair, or new agent trial. Historical v3/v4 data remain in the
appendices with their original cohort identities.

\begin{table}[t]\centering\small
\caption{New v5 synthetic study. Reduction is selected-layout held-out padded cost relative to the original-order balanced incumbent; positive is better. Fit includes affinity/order learning and catalogue construction, not validation scoring. Three seeds per row; ranges are descriptive.}\label{tab:v5synthetic}
\begin{tabularx}{\linewidth}{@{}l l r r r l@{}}\toprule
Checks & Family & Mean (\%) & Range (\%) & Fit median (s) & Changed\\\midrule
128 & cluster & 23.43 & 22.40--24.45 & 0.018 & 3/3\\
128 & independent & 0.00 & 0.00--0.00 & 0.016 & 0/3\\
128 & global & 0.00 & 0.00--0.00 & 0.032 & 0/3\\
512 & cluster & 10.75 & 10.06--11.13 & 0.523 & 3/3\\
512 & independent & 0.07 & 0.00--0.20 & 0.548 & 1/3\\
512 & global & 0.00 & 0.00--0.00 & 0.645 & 0/3\\
1{,}024 & cluster & 5.56 & 5.16--5.80 & 2.991 & 3/3\\
1{,}024 & independent & 0.00 & 0.00--0.00 & 3.086 & 1/3\\
1{,}024 & global & 0.00 & 0.00--0.00 & 3.462 & 0/3\\
\bottomrule\end{tabularx}\end{table}

\subsection{Practical default without exact or block DP}\label{sec:practical}
On clustered completions, mean held-out reductions are 23.43\%, 10.75\% and
5.56\% as registry size increases (Table~\ref{tab:v5synthetic}). Global updates
give no reduction; independent updates give zero except for one of the three
512-check seeds, whose small gain produces a 0.07\% mean. These controls do not
support universal value from ordering. The selected method includes no DP, and a
regression test makes any call to the exact fitting helper fail on this path.

The full catalogue fit median at 1{,}024 checks remains about 3 seconds; order
learning is not free. The fixed construction can reuse a previously approved
order in linear work, but no linear-time claim is made for the full learned
pipeline. The smaller percentage at larger registries is consistent with
substantial layout-independent cold construction in this benchmark; it should not
be interpreted as a learned scaling law.

An initial execution exceeded its 120-second tool limit after saving 26 complete
records. No process remained; explicit resume skipped those exact records and
completed the remaining frozen seed without changing data or selection. All 27
observations and the interruption record are retained.

\subsection{Source-evidence replays: simplicity is not byte optimality}\label{sec:replays}
Learned balanced layouts reduce whole-epoch compressed bytes by 41.84\% for
NetworkX, 4.36\% for Toolz and 18.53\% for Future Code (Table~\ref{tab:replay}).
These are not the larger v4 exact-layout savings relabeled as results of the new
default. The balanced representation remains larger than learned exact in every
project here. Exact fitting may therefore be appropriate when amortized delivery
justifies its additional compilation; the practical default trades that potential
gain for simpler fitting.

Across the twelve cases, batch-compressed application bytes are 416{,}046 for
fixed8, 299{,}081 for learned balanced and 278{,}831 for the re-serialized
learned exact layout. Semantic counts, verdicts and local authorization agree
across all 36 replays. Their provenance binds the archived check records by
digest. Serialization includes newly bound layout/run metadata; numerical bytes
need not exactly reproduce a prior archive's compressibility. The independent
audit checks actual emitted bytes rather than substituting earlier totals.

\begin{table}[t]\centering\small
\caption{New serialization of twelve archived source candidates, four per project. The first three columns compare learned-order balanced8 with original-order fixed8. The last column compares learned balanced with learned exact after epoch compression: positive is a penalty. All include the codec's application framing.}\label{tab:replay}
\begin{tabularx}{\linewidth}{@{}l r r r r@{}}\toprule
Project & Padded (\%) & Canonical (\%) & Epoch (\%) & vs.\ exact penalty (\%)\\\midrule
networkx & 48.63 & 48.74 & 41.84 & 11.68\\
toolz & 8.78 & 8.70 & 4.36 & 1.52\\
future & 29.12 & 29.23 & 18.53 & 6.67\\
\bottomrule\end{tabularx}\end{table}

Root-status text is identical in all twelve triples and occupies 593--594 UTF-8
bytes. Exhaustive diagnostic pages total 148 for fixed8, 95 for learned balanced
and 87 for learned exact. The new default therefore uses 35.81\% fewer pages than
fixed8, not v4's 41.22\% exact-layout reduction. The largest emitted page is
16{,}364 bytes under a 16{,}384-byte limit. Reading every page is a specified
exhaustive diagnostic service, not observed LLM behavior. There is no measured
token reduction; root token equality follows from identical text for a fixed
tokenizer, but no numerical BPE counts are asserted.

\subsection{Codec-matched deployment scenarios, not measured speedups}\label{sec:scenarios}
The measured fixed-to-balanced compressed byte difference is 116{,}965 per
twelve-candidate cohort. Equation~\eqref{eq:netbenefit} gives 1.784744 seconds of
gross delivery saving per cohort copy at an assumed 64\,KiB/s goodput, with
$E=0$. This is arithmetic from replayed byte sizes, not the older 2.21-second
measured loopback saving of a different layout.

The 5-second column (Table~\ref{tab:qmin}) is a deliberately declared sensitivity
scenario, not an estimate transferred from the historical harness. At 64\,KiB/s
and one export per candidate, the assumed 60 seconds of cohort overhead
overwhelms the approximately 1.785-second delivery saving. Repeated delivery
might change the equation, but the application must actually require it. With new
candidates on each repetition, preparation recurs and cannot be amortized as the
same $C$. All 27 scenarios (three codecs, three rates, three overhead assumptions)
are available, including negative outcomes. The planner rejects mismatched
services and missing measurements rather than inventing a benefit.

\begin{table}[t]\centering\small
\caption{Conditional minimum profitable copies of the same twelve-candidate cohort with whole-epoch compression. $E=0$ is an assumption. The columns assume extra cost per candidate, paid once; cohort costs are 0.6, 6 and 60 seconds. These are not measured strict-wrapper overheads.}\label{tab:qmin}
\begin{tabularx}{\linewidth}{@{}l r r r@{}}\toprule
Assumed goodput & 0.05 s/candidate & 0.5 s/candidate & 5 s/candidate\\\midrule
8\,KiB/s & 1 & 1 & 5\\
64\,KiB/s & 1 & 4 & 34\\
1\,MiB/s & 6 & 54 & 538\\
\bottomrule\end{tabularx}\end{table}

\subsection{Historical controls remain controls}\label{sec:controls}
The earlier 288 loopback transfers belong to v4, use application pacing rather
than real WAN conditions, and are not rerun or relabeled as v5 trials. For twelve
compressed fixed/exact exports at 64\,KiB/s, their measured totals were 6.624 and
4.413 seconds. The larger uncompressed totals were 33.499 and 18.954 seconds.
Compression of fixed8 was more impactful than selecting an uncompressed learned
layout. None of those durations is an IP/TLS throughput measurement or predicts a
loss-prone path.

The original 39.69/39.77-second checker totals and 12.98/18.11/18.12-second
median direct/fixed/learned harness durations remain negative controls. They came
from different source/harness fixtures, so this paper does not subtract their
times across cohorts. The historical window-controller and block-compiler
evidence is retained unchanged. Block DP's 1.54--5.34\% penalty to same-order
balanced8 is the reason it is now an ablation, not a contribution advertised as
deployment acceleration.

\subsection{Unmeasured endpoints and tested interfaces}\label{sec:unmeasured}
Both requested BPE encodings, \texttt{cl100k\_base} and
\texttt{o200k\_base}, remain \UNKNOWN{}: the software was absent,
package/vocabulary fetches failed, and a connector-retrieved build artifact was
expired. No substitute tokenizer, model alias or byte division was used. The
measurement tool records exact content-only counts when a real encoding is
available; its unit tests use explicitly labeled toy counters only to validate
admission logic.

No two authorized remote measurement hosts or designated live model endpoint were
available. Physical-WAN and hosted-LLM trial counts are zero. The new TLS
utility's local compatibility tests check transport integrity, certificate
rejection and bounded framing; they do not measure WAN jitter/loss or close the
agent-evaluation gap. A reproducible deployment protocol is supplied, rather than
a claimed remote experiment.

\subsection{Software and independent audit}\label{sec:audit}
The prototype suite now spans both protocol generations of the
retirement study: the v1 engine's 10 pinned invariants
(determinism, summary/reset failure, full-replay coverage, flat
steel-man accounting, closure outcomes) and the v2 engine's 12
(factorial coverage, scheduler byte-neutrality, phase-sum identity,
audit cleanliness, closure directions, long-horizon persistence,
bit-identical fast layouts). The v2 suite caught a real defect during
development---the audit's ground-truth pending set deduplicated
obligation openings by first rather than latest opening---which is
disclosed in Section~\ref{sec:retire-results}; the frozen cohort was
regenerated before any v2 manuscript text was written. The earlier
serialization prototype passes 329 tests, including 76 new cases; the existing
version-equality assertion is updated to require 5.0.0 consistently in runtime and
package metadata. A duplicate test-module basename and the initially inconsistent
version were corrected without weakening behavior checks; failure logs are
retained. New tests cover no-DP compilation, service mismatch, missing costs,
exact break-even, whole-text token admission, TLS trust failures and deliberate
data tampering.

A separate auditor imports no producer optimizer, selector or codec. It
reconstructs full-initialization and wave costs for all 27 workloads, checks 36
serialized replays against original evidence digests and verdicts, reconstructs
all codec byte sizes and diagnostic pages, and independently checks 27 planning
scenarios and 131 summary comparisons. These checks validate recorded arithmetic
and provenance, not third-party security certification. Release verification
separately records extracted-source tests, the unchanged Future Code suite,
installed-source hashes and archive integrity.

\section{Deployment Recommendation, Limitations, and the Pre-Registered Live Cohort}\label{sec:discussion}
\subsection{Recommendations from the simulator, and the protocol-only tier}\label{sec:deploy}
The recommendations below have two tiers: what this paper's simulator
measures, and what remains protocol-only.

\paragraph{Mechanism-level (measured in this paper's simulator).}
Where a service already requires protected verification evidence and a
constrained channel, the serialized-evidence component of
Section~\ref{sec:substrate} applies as before. The
retirement study adds a sharper pattern for context policy: in the
simulator, the policy that first converts history into current
verified bindings, a typed pending-obligation set, and dependency
interfaces passes the acceptance contract at every scale tested,
while policies that rely on lossy summaries of obligations, or that
revert values to beliefs older than the current working window,
measurably fail (Section~\ref{sec:retire-results}). The
conversion is the pattern the simulator supports; whether it is worth
adopting depends on the
cost regime and on the retrieval layer. It buys \RetwoOpsReduction{}
fewer retrieval round trips at equal correctness, and by the closed
form of Proposition~\ref{prop:breakeven} the simulator's gate flips
at a computable round-trip price: given means
$(B_f,B_m,o_f,o_m)$ measured in a deployment's own pilot and that
deployment's price $c$, the formula's threshold $c_{0.8}$ states
where the frozen 20\% gate would flip. In the simulator,
$c_{0.8}\approx1{,}364$\,B pooled (frontier's 8.0\% byte advantage
does not by itself clear the 20\% gate; the frozen 4{,}096\,B
round-trip price does). These are simulator thresholds under the
frozen accounting, not validated deployment guidance: they identify
what a deployment must measure---its own byte costs, operation mix,
and round-trip price---rather than prescribing adoption.
Second, and now stated because the exploratory analyses forced it: a
64-gate bounded-recency window (\texttt{bounded-hier}) is cheaper
than frontier below a break-even price of
\ExpBhBreakEvenPooled{} pooled (2{,}148\,B independent; never in
global) at equal acceptance---but \ExpBhDueExposed{} of its due
obligations and \ExpBhFootprintExposed{} of its footprint reads sit
outside its retention, so its safety delegates to the retrieval
layer. The simulator-conditional reading is therefore a function of
two observable quantities, not one: with reliable, cheap retrieval and
round-trip price below the break-even, the bounded window is the
cheaper arm in the simulator; with expensive or unreliable retrieval,
the conversion heuristic is the arm whose model-internal guarantees
do not depend on the retrieval layer beyond the
fetch itself. Which of these a real deployment should prefer is
exactly what the live cohort measures; the simulator brackets the
possibilities rather than deciding between them. The real-trace replay
(Section~\ref{sec:retire-exploratory}) brackets what this looks like
on non-synthetic change patterns. Between grouped
certificate layouts and flat delta ledgers, select by update density:
the measured-density rule matched the per-cell oracle on every cell of
the frozen cohort (\RetwoDensityRegret{} regret); dense updates favor
the tree, sparse deltas favor the flat ledger, and the density signal
is observable on a training half of any real stream. The
demand-driven adaptive variant is not promoted: its pre-registered
gate failed in the simulator
($\RetwoAdaptiveReduction{}$ versus plain frontier), and adoption
decisions belong to the live cohort, not to the simulator alone.

\paragraph{Live-agent cohort (pre-registered, not run; highest-priority next experiment).}\label{sec:live-protocol}
The mechanism study predicts, but cannot confirm, that avoided
retrieval operations translate into live savings. Executing this
cohort is the single most informative next step the study permits
itself: it is the only pre-registered instrument that can move
real-world agent relevance from moderate to measured, and every
limitation this section records---BPE token counts, provider billing,
prompt-semantics effects, and live tool-call latencies---is either an
outcome it measures or a risk it prices. The confirmatory
cohort is therefore pre-registered with operations as the primary cost
measure:
\begin{enumerate}[leftmargin=*,itemsep=1pt]
\item \textbf{Workload.} Multi-stage changes on a fixed public issue
benchmark (SWE-bench-class), screened for dependency coupling; the
same frozen independent acceptance contract for all arms; equal total
budgets, sandbox, and protected state; every failed branch in the
numerator.
\item \textbf{Arms.} The arms the simulator separates:
observation masking; summary/reset; verified-frontier retirement; and
the density-selected evidence layout (flat delta ledger versus
grouped tree, chosen per the measured update density of the pilot).
The demand-driven adaptive variant is excluded because its promotion
gate failed in the simulator; including it would re-litigate a
pre-registered decision.
\item \textbf{Primary outcomes.} Tool-call round trips per accepted
project (the pre-registered live-cost measure, from the mechanism
result); cost per accepted project (all provider metering preserved,
unknowns never zero-imputed); acceptance rate; deadline success with
censored timeouts. Pre-specified noninferiority margin on acceptance:
no arm may trail masking by more than 5\% absolute.
\item \textbf{Closure rules.} \emph{Null closure}: if frontier
retirement does not reduce tool-call operations per accepted project
by at least 20\% versus masking at noninferior acceptance, the
live-cohort hypothesis is rejected and retirement is not promoted.
\emph{HACT closure}: if the density-selected layout does not beat
flat-delta evidence by at least 5\% on measured evidence bytes at
equal correctness, the tree is not deployed.
\end{enumerate}
Two explicit closure rules prevent indefinite iteration; either rule
terminates the program with a written result. This protocol is frozen
here, before any live data collection.

The live layer also prices risks the mechanism study cannot express
and must not silently assume away. A hosted model does not read
idealized state capsules: structured JSON or capsule formats can
induce attention degradation, prompt distraction, or formatting drift,
so the cohort scores each arm's full delivered prompts, not just their
sizes. Real summarization does not follow a synthetic four-episode
boundary, so the summary/reset arm keeps its own provider-side
implementation with the acceptance contract as the only fixed point.
And live round trips are not a static additive constant: network
jitter, rate limiting, and variable generation latency replace the
frozen $4{,}096$\,B accounting constant with the cohort's measured
per-tool-call durations, which is why tool-call counts, not byte
totals, remain the pre-registered primary outcome.

\paragraph{What remains unmeasured.}
No hosted model was called for any number in this paper; BPE
tokenization, physical-WAN behavior, and provider billing remain
unknown and are never zero-imputed. Two tokenization caveats deserve
emphasis because the paper's savings are reported in bytes. First,
both requested BPE encodings (\texttt{cl100k\_base} and
\texttt{o200k\_base}) remain \UNKNOWN{}
(Section~\ref{sec:unmeasured}); byte savings therefore do not
necessarily convert to proportional BPE token savings, and no
byte-to-token division is asserted anywhere. Second, even a measured
conversion would be model-specific: tokenization is not a
length-preserving map, so the live cohort records actual tokenizer
output per arm rather than converting the byte result. The mechanism cohort's cost model
is a frozen accounting of a frozen model of information flow; real
payloads are not lognormal, and real summaries may retain obligations
better or worse than the seeded boundary survival. The round-trip
regime's $4{,}096$\,B constant is an accounting choice bracketed by a
sensitivity range, not a measured latency. The tool-call prediction
is the only bridge from the mechanism result to live cost, and it is
pre-registered as such rather than assumed.

\paragraph{Trust and scope.}
Future Code remains a single-host supervisor. The guard and TLS
utility are not a distributed consensus protocol or fault-tolerant
deployment. A passing declared suite cannot establish absence of
arbitrary defects. Authorization requires an uncompromised checker,
correct protected reference data, immutable candidate identity, and
local monotone state; Theorem~\ref{thm:soundness} stands as before.
TLS authenticates the configured server, not the checker's
truthfulness. The receiver is a bounded measurement utility behind an
operator-controlled firewall, not a production ingestion service.

\appendix

\section{Audit and Reproduction}\label{app:repro}
\subsection{Exact data generation}
The synthetic generator uses NumPy's \texttt{default\_rng} initialized with a
\texttt{SeedSequence} whose entries are \texttt{[260918, n, seed, stream]}.
Stream zero produces training episodes and stream one held-out episodes. Group
construction uses \texttt{SeedSequence([7001,seed,n])}. Pilot seeds 1000 and 1001
are excluded from all main tables. Scaling sizes other than 128 have their own
saved records; the four reused 128-gate seeds are clearly identified in the
analysis.

The v5 frozen contract fixes sizes 128, 512 and 1{,}024; families clustered,
independent and global; and seeds 55011--55013, with 128 training, 128 validation
and 256 held-out epochs per workload. The executable fixture uses
\texttt{SeedSequence([83012,seed,stream])}. Interrupted fixture runs append
completed snapshot records to a journal and replay those records when resuming;
they do not reclassify incomplete work as a completed seed.

The main commands, from the extracted research root, are:
\begin{verbatim}
python -m pytest -q
python -m experiments.synthetic --family clustered --n 128 --seeds 20
python -m experiments.executable --seeds 6 --out results/executable_fresh
python -m experiments.sensitivity
python -m experiments.aggregate
python -m experiments.validate_results
\end{verbatim}
Nothing in these commands contacts an LLM endpoint. The artifact separates raw
results from regenerated tables.

\subsection{Software and data identities}
Archived runs used Python 3.13.5 on Linux x86\_64; exact installed package and
platform records accompany emitted reports. Compilers were measured with
\texttt{OPENBLAS\_NUM\_THREADS=1}. No accelerator is required. Vendored upstream
sources and tests are retained with their notices. Every fresh mutation is checked
against its archived base-source hash, replacement span, and resulting source
hash. This binds a result to a specific intervention, not a natural historical bug
distribution.

\section{Certificate Encoding and Trust-Boundary Details}\label{app:encoding}
Each capsule contains \texttt{range}, \texttt{counts}, \texttt{binding},
\texttt{evidence}, and \texttt{registry}. Hash fields are 64-character lowercase
SHA-256 hexadecimal encodings. Canonical JSON uses sorted keys, ASCII escapes, and
compact separators. In a packet, the header occupies 512 bytes including its
newline; each child card occupies 320 bytes including its newline. Padding makes
byte accounting exact. Fixed padding is not asserted to be a token-optimal
production encoding; a binary or tokenizer-aware encoding would define a new price
function and require re-evaluation.

The trusted verifier API rejects mismatched bindings, invalid verdicts, and
malformed trace hashes. It is intentionally not exposed as a tool through which an
untrusted LLM can certify its own work. Returned display capsules are compared
with canonical trusted records rather than accepted because their self-reported
hashes appear valid. This local design demonstrates record integrity and freshness
under its trust assumptions; distributed authenticated transport and credential
isolation remain host-runtime responsibilities.

The optimizer's $O(bn^3)$ claim concerns planning. The reference kernel
prioritizes clarity and scans interval generation stamps; it also hashes tree
descriptions. We do not infer an asymptotically optimal runtime-update
implementation from the communication bound. Large-scale implementations could use
range-maximum indexes and cached topology digests, but those optimizations are
outside the evaluated prototype.

\section{Optional Fixed-Share Analysis: Not the Deployment Recommendation}\label{app:fixedshare}
This appendix retains the former tracking theorem and its assumptions. It is an
optional comparator, not the practical selector. On the measured streams it loses
to the deterministic priced-window rule in three regimes. We have not searched for
a favorable adversarial stream after seeing those results, and make no claim that a
tracking bound implies an empirical advantage.

Let a declared bound $R$ cover each round's loss spread,
$\max_i\ell_t(i)-\min_i\ell_t(i)\le R$. Reject a measurement exceeding it instead
of clipping the cost. Define $z_{t,i}=(\ell_t(i)-\min_j\ell_t(j))/R$. Starting at
$p_1(i)=1/K$, update
\begin{align}q_{t+1}(i)&=\frac{p_t(i)\exp(-\eta z_{t,i})}{\sum_j p_t(j)\exp(-\eta z_{t,j})},\label{eq:fixedshare}\\
p_{t+1}(i)&=(1-\alpha)q_{t+1}(i)+\alpha/K,\qquad 0<\eta\le1,\ 0\le\alpha<1.\label{eq:mixing}\end{align}
Independent draws from each new distribution can waste migrations. Instead,
conditional on $A_t=i$, retain $i$ with probability
$\min\{1,p_{t+1}(i)/p_t(i)\}$. If it is not retained, sample from the normalized
positive residual $(p_{t+1}-p_t)^+$. This maximal coupling has marginal $p_{t+1}$
and switch probability $\mathrm{TV}(p_t,p_{t+1})$. Established dependent-sampling
methods motivate this choice~\cite{geulen}; we do not claim it as a new sampling
algorithm.

\begin{theorem}[Catalogue tracking with migration]\label{thm:tracking}
Assume a frozen catalogue, a fixed exogenous sequence of full-information loss
vectors, spread at most $R$, and $M_{ij}\le M$. For any comparator path $u_{1:T}$
with $s$ changes and $0<\alpha<1$, the coupled policy satisfies
\begin{equation}\E[L_T]-\sum_t\ell_t(u_t)\le
R\left(\frac{\log K+s\log(K/\alpha)+(T-1-s)\log(1/(1-\alpha))}{\eta}+\frac{\eta T}{8}\right)
+M(T-1)(\eta+\alpha).\label{eq:tracking}\end{equation}
For $\alpha=0$ the stationary-comparator bound has numerator $\log K$ and applies
only to $s=0$.
\end{theorem}
\begin{proof}
The fixed-share update is exponential weighting over paths with an initial uniform
prior, self-transition mass at least $1-\alpha$, and transition mass at least
$\alpha/K$. The negative log prior of a path with $s$ switches is bounded by the
numerator. The bounded exponential-moment inequality gives expected normalized
service regret at most that numerator divided by $\eta$, plus $\eta T/8$. Since
losses are exogenous, the weights are independent of the sampled action path, and
the coupling preserves each unconditional marginal. Exponential reweighting
changes total variation by at most $\eta$; uniform mixing adds at most $\alpha$.
Thus the expected switch count is at most $(T-1)(\eta+\alpha)$.
Multiplication by $M$ and addition of service regret proves the bound. \qed
\end{proof}
The bound is not a per-run improvement guarantee. Choosing fixed $\eta=.5$ and
$\alpha=.01$, as in our frozen experiment, need not make this bound sublinear with
increasing $T$.

\section{Block-Backbone Compiler: Retained Ablation}\label{app:block}
The exact compiler's cubic cost makes frequent full recompilation inappropriate.
For a supplied display order, divide the leaves into consecutive blocks of at most
$s$ leaves, and fix a compact $b$-ary macro-tree on those blocks. Optimize each
block's microtree using its restricted completion waves and strict-epoch
initialization cost. Replace each macro leaf by its optimized microtree. No
evaluation waves enter this process.

The proposition below is exact only under its stated restrictions, and one
restriction deserves emphasis because it is easy to misread as a general
claim: \emph{every update wave that intersects a macro node's span must
cover whole blocks within it}. If a wave can hit a partial block---some
but not all of the leaves a macro node covers---then that node's
activation weight varies within the class and the objective no longer
separates into per-block subproblems; the $O(bns^2)$ claim then fails
and the decomposition is only a heuristic. The frozen streams of this
paper never emit partial-block waves (a wave's gate set is a union of
whole invalidation sets, and the measured waves intersect macro spans
block-completely), which is why the decomposition is exact on them; a
deployment with partial-block waves must either re-block or fall back
to the full exact compiler.

\begin{proposition}[Exactness in the block-backbone class]\label{prop:block}
With display order, block boundaries and macro-tree fixed, independent exact
optimization of the blocks minimizes the training objective over every layout in
that restricted class. The dynamic-programming work is $O(bns^2)$ with $O(bs^2)$
reusable DP workspace, excluding order learning and wave preparation.
\end{proposition}
\begin{proof}
Every macro node covers a fixed union of complete blocks, so its activation weight
and packet price are constant across the class. All variable costs lie inside
disjoint microtrees. The objective is a constant plus the sum of their independent
objectives. Conditional exactness therefore gives the minimizing member. There are
at most $\lceil n/s\rceil$ subproblems, each of size at most $s$, yielding the
stated DP bound. Reusing the subproblem arrays yields the workspace bound. \qed
\end{proof}
For the reference implementation, $M$ recorded waves and $L$ total wave
memberships add filtering and weight-preparation work of
$O(\lceil n/s\rceil L+Mn s)$, and storage for the input, output and local $M\times
s$ arrays. This is not a linear-time claim for the entire pipeline. No bounded
approximation ratio against the unrestricted exact tree is asserted; held-out cost
can deteriorate. The historical experiment froze $s=32$ before measurement. Its
negative result----block DP lost to same-order balanced fixed8 in held-out
bytes----motivates excluding this compiler from the practical path.

\section{Information Separations}\label{app:sep}
These results concern the information available to a selector. They are not
hardness results for all adaptive tree algorithms, and they do not imply that a
complete hypergraph must be recovered: interval hit statistics are sufficient for
a fixed order and objective.

\subsection{Marginals do not determine a co-invalidation optimum}
Consider four ordered leaves, binary trees, and a constant internal packet price
$c>0$. Let $\mathcal D_A$ choose $\{0,1\}$ or $\{2,3\}$ equiprobably and
$\mathcal D_B$ choose $\{0,3\}$ or $\{1,2\}$ equiprobably. All leaf marginal
activation probabilities equal $1/2$ in both distributions. The root and every
three-leaf interval are always activated. The five tree shapes fall into the
following classes:
\begin{center}
\begin{tabular}{lc c}\toprule
Ordered tree & $J_A/c$ & $J_B/c$\\\midrule
$((0,1),(2,3))$ & 2 & 3\\
$(0,((1,2),3))$, $((0,(1,2)),3)$ & 3 & 2.5\\
$(0,(1,(2,3)))$, $(((0,1),2),3)$ & 2.5 & 3\\
\bottomrule\end{tabular}
\end{center}
Optimal costs are $2c$ and $2.5c$. A deterministic marginal-only selector receives
identical input and must choose the same class, incurring at least $c/2$ regret on
one distribution. The last class is weakly dominated by the balanced tree. If a
randomized selector chooses the balanced tree with probability $q$, its regrets
are at least $(1-q)c$ and $qc/2$. Their maximum is minimized at $q=2/3$ and equals
$c/3$, which is attainable for this instance. A fixed initialization price adds the
same $3c$ to all binary shapes and does not alter the separation.

\subsection{Joint invalidation does not determine progressive layout}
Every epoch now invalidates all four checks, and execution order is $(0,1,2,3)$
with singleton completion waves. In workload $A$, check 0 fails, so only its wave
is executed. In workload $B$, check 3 fails, so all four checks execute. Initial
invalidation is identical, even jointly. Every binary tree has three internal
nodes; its cost in $A$ is $(3+\mathrm{depth}(0))c$ and in $B$ is
$(3+\sum_i\mathrm{depth}(i))c$. The root-0 comb minimizes $A$ at $4c$ but costs
$12c$ in $B$, whereas the balanced tree minimizes $B$ at $11c$ but costs $5c$ in
$A$. An invalidation-only randomized selector using a root-0 comb with probability
$p$ incurs regrets at least $(1-p)c$ and $pc$, whose maximum is at least $c/2$. A
deterministic selector incurs at least $c$. This proves a separation between
stale-set information and schedule-induced completion information; it does not
prohibit a selector that observes both.

\section{Exact Conditional Tree Synthesis: Reference Dynamic Program}\label{app:exactdp}
The practical default of Section~\ref{sec:compiler} never calls this dynamic
program; it is retained as the exact reference the catalogue construction is
compared against, and it is invoked only by the historical research compilers.

Let $D(i,j)$ be the optimal subtree cost and $F_k(i,j)$ the optimal cost of an
ordered forest with $k$ nonempty children. With infeasible states equal to
$+\infty$,
\begin{align}
D(i,i)&=0,\label{eq:Dbase}\\
F_1(i,j)&=D(i,j),\label{eq:F1}\\
F_k(i,j)&=\min_{i+k-2\le u<j}\bigl\{F_{k-1}(i,u)+D(u+1,j)\bigr\},\label{eq:Forest}\\
D(i,j)&=\min_{2\le k\le\min(b,\,j-i+1)}\bigl\{r_\pi(i,j)c(k)+F_k(i,j)\bigr\}.\label{eq:D}
\end{align}

\begin{theorem}[Exactness conditional on order]\label{thm:exact}
For a fixed $\pi$, nonnegative interval weights, and the packet model in
Equation~\eqref{eq:price}, Equation~\eqref{eq:D} returns a minimizing feasible
tree in $O(bn^3)$ time and $O(bn^2)$ space.
\end{theorem}
\begin{proof}
Every feasible root has an arity $k$ in the stated range, so the outer
minimization ranges over exactly the feasible root arities. Each candidate tree
splits into a root and an ordered forest of $k$ subtrees whose positional
intervals partition $[i,j]$; conversely every such split yields a feasible tree.
The recurrence over $F_k$ explores exactly these splits, and the induction on
interval length proves exactness.
There are $O(bn^2)$ states, each scanning at most $n$ splits. The apparent
$F_1=D$ dependency is used only on shorter intervals when constructing a forest
with at least two components. \qed
\end{proof}
This is a specialization of classical alphabetic-tree
optimization~\cite{alphabetic}. The historical three-proposal reference
multiplies compilation time by a constant, not by $n!$. Dense spectral
computation can itself cost $O(n^3)$; the reference clustering implementation
uses $O(n^2\log n)$ heap operations and $O(n^2)$ storage.

\section{Cost, Statistical, and Online Proof Details}\label{app:proofs}
\subsection{Full-refresh lower bound and global negative control}
For an $n$-leaf tree with $m$ internal nodes, the number of edges is $n+m-1$, so
$\sum_v d_v=n+m-1$. Therefore
\begin{equation}W(L)=S(n-1)+(H+S)m\ge S(n-1)+(H+S)\left\lceil\frac{n-1}{b-1}\right\rceil.\label{eq:lb}\end{equation}
The inequality follows from $n-1=\sum_v(d_v-1)\le m(b-1)$. At $n=64$, $b=8$, a
complete eight-child construction attains the bound with $m=9$ and $W=27{,}648$
bytes. Thus all selected global-control layouts can tie an analytically minimal
full refresh. Strict initialization followed by one global completion wave costs
$2W$. This explains the exact negative control without assuming every balanced tree
at every $n$ attains the lower bound.

\subsection{Uniform statistical control at a fixed order}
For one activation wave per episode and a fixed registry order, let
$M_n=n(n+1)/2$ and $\epsilon=\sqrt{\log(2M_n/\delta)/(2m)}$. The interval hit
indicators are bounded independent variables across independent episodes, though
checks within an episode may be arbitrarily dependent. Hoeffding plus a union bound
gives simultaneous interval-frequency error at most $\epsilon$. For every feasible
tree, $|\widehat J(T)-J(T)|\le W(T)\epsilon$. Inserting the empirical minimizer
gives regret at most $[W(\widehat T)+W(T^*)]\epsilon\le 2(n-1)C\epsilon$. This is
the original HACT bound, retained for attribution. It is not automatically valid
uniformly over arbitrary data-selected permutations; the separate held-out
catalogue bound in Proposition~\ref{prop:valid} addresses that selection step.

For strict epochs with at most $q$ completion waves, a loose range is
$B=(q+1)(n-1)C$. Under the evaluated sequential eight-record flush rule with at most
one final record per obligation, $q\le\lceil n/8\rceil$. A narrower observed range
is not automatically a proved bound for an unobserved deployment distribution.

\subsection{Maximal coupling and horizon dependence}
For consecutive distributions $p,p'$, the probability mass that remains at index
$i$ under the coupling is $\min(p_i,p'_i)$. The remaining mass is distributed as
$(p'-p)^+$, so the new marginal is exactly $p'$ and $\Pr(A_{t+1}\ne A_t)=\mathrm{TV}(p,p')$.
For exponential reweighting by $a_i=\exp(-\eta z_i)$ with $z_i\in[0,1]$, the
normalizer is at most one, hence $q_i\ge e^{-\eta}p_i$. Thus
$\sum_i\min(p_i,q_i)\ge e^{-\eta}$ and $\mathrm{TV}(p,q)\le1-e^{-\eta}\le\eta$.
Mixing $q$ with the uniform distribution by weight $\alpha$ changes total variation
by at most $\alpha$. The triangle inequality gives the switch bound in
Theorem~\ref{thm:tracking}.

For $K\ge2$, $s=0$, and $\alpha=0$, write $B_T=RT/8+M(T-1)$. The bound is
$R\log K/\eta+\eta B_T$. When $\sqrt{R\log K/B_T}\le1$, choosing that learning rate
gives at most $2\sqrt{R\log K\,B_T}$ expected regret, sublinear for fixed
$R,M,K$. This is an established style of switching-cost bound, not a new rate. A
single empirical switch with estimated future saving $d$ and anticipated horizon
$h$ satisfies the familiar break-even inequality $hd>M_{ij}$. This is arithmetic,
not a drift theorem: savings may reverse before $h$ events.

\section{Evidence-to-Claim Register}\label{app:register}
\begin{center}\small
\begin{tabularx}{\linewidth}{@{}p{.30\linewidth}Xp{.22\linewidth}@{}}\toprule
Claim & Evidence & Scope boundary\\\midrule
Exact ordered optimum & Theorem~\ref{thm:exact}; exhaustive enumeration tests & Fixed order and price model\\
Validation after order & Proposition~\ref{prop:valid}; Hoeffding bound & Fixed catalogue, no shift\\
Publication soundness & Theorem~\ref{thm:soundness}; conditional proof & Trusted checker, protected state\\
Practical byte reduction & Table~\ref{tab:v5synthetic}; frozen seeds & Clustered workloads, 3 seeds\\
Source replay reductions & Table~\ref{tab:replay}; 36 replays & Diagnostic, not fresh checks\\
Root status identical & Equation~\eqref{eq:rootid}; 593--594 bytes & Fixed tokenizer, layout-independent\\
Break-even accounting & Equations~\eqref{eq:netbenefit}--\eqref{eq:qmin} & Declared $C,E,R$, codec-matched\\
BPE tokenization gap & \texttt{cl100k\_base}/\texttt{o200k\_base} remain \UNKNOWN{}; admission tool supplied & Vocabulary retrieval failed; no billing claim\\
Context-retirement mechanism study & Pre-registered, factorial, executed: v2 = 1{,}176 runs + 63 long-horizon, seven retention levels $\times$ scheduler factor, five closures (two proposals rejected by their gates) & Measured at mechanism level; live cohort still protocol-only\\
\bottomrule\end{tabularx}\end{center}

\section{Retirement-Study Protocol Deviations and Engine History}\label{app:retire-deviations}
Two protocol generations were run. Generation v1 (file-frozen
2026-09-23; \texttt{docs/\allowbreak protocols/\allowbreak CONTEXT\_\allowbreak RETIREMENT\_\allowbreak PROTOCOL\_\allowbreak v1.md}
with amendments v1.1--v1.4) was written before cohort data generation, and
each amendment was written before the run it governs. The engine went
through four substantive revisions during that pre-data window; the
final cohort was generated once, by the amended engine, and never
regenerated after any manuscript text was written. For auditability we
disclose the full history, including the errors that motivated each
amendment:

\begin{enumerate}[leftmargin=*,itemsep=2pt]
\item \textbf{Amendment v1.1 (pre-data).} The HACT workload generators
draw invalidation sets i.i.d.; retirement is temporal, so persistence
transitions (0.6) and fixed obligation windows (six episodes) were
added. A single-stream smoke test of the original design exposed
nominal (name-based) obligation retention; the amendment made
retention emergent from each policy's mechanism.
\item \textbf{Amendment v1.2 (pre-data).} A 630-run smoke test to a
temporary directory (results void, never cited) exposed three
under-specifications: hand-rolled evidence accounting, summary/reset
modeled as sleeping between resets rather than attending, and a
frontier capsule that did not scale with contents. All were pinned
before the final run.
\item \textbf{Amendment v1.3 (pre-data).} A single-stream check
showed facts carrying no payload, making replay artificially cheap;
payloads (lognormal, median 128\,B), payload-charged retrieval, and
interface-only capsules were pinned.
\item \textbf{Amendment v1.4 (pre-data).} Final verification exposed
missing dependency footprints (the review's exported interfaces), an
inconsistent summary/reset failure model, and the need for retrieval
operations as a co-primary outcome. The evidence ablation was refined
to a held-out half with a learned layout order and a steel-manned
flat ledger taking the better of full and delta exports.
\item \textbf{Post-freeze fixes.} A flat-ledger registry-size
off-by-one was corrected after the first final-format cohort; the
cohort was regenerated once, before any manuscript table was
generated from it. One obligation-resolution accounting error
(retained-value reads not crediting resolution) was fixed during
pre-data engine verification; all obligation flows are now covered by
the test suite (\texttt{tests/retirement}), which also pins
determinism, the closure outcomes, and the per-family evidence
directions.
\end{enumerate}

No amendment was written after the final cohort was generated. The
v1 cohort, closures, and summary JSON are reproduced by
\texttt{bun src/harness/retirement/cli.ts} in under two minutes and
are byte-identical across runs; the v2 cohort, closures, and tables
are reproduced by \texttt{bun src/harness/retirement/run2.ts} and
\texttt{bun src/harness/retirement/tex2.ts} with the same property.
Both engines, both protocols, and both frozen cohorts are committed;
v2's protocol and engine-notes commits precede its data commit, which
precedes the first v2 manuscript-text commit.


\section{Generation-v1 Record (Complete)}\label{app:retire-v1}
The v1 study is superseded by the v2 factorial
(Section~\ref{sec:retire}) but is retained in full for provenance: v2's
design was chosen to repair v1's weaknesses, and v1's closures
motivated v2's frozen rules. The v1 protocol was frozen in files on
2026-09-23, before its cohort ran; it was first committed to git on 2026-09-24, after the v2 protocol commit
(333d495). v1's freeze is
therefore file-attested, and the ordering is disclosed rather than
smoothed.

\paragraph{Design.} Ten arms over $3\times3\times7=63$ cells (630
runs, 40 episodes each): full-replay, masking, summary-reset,
bounded-hier, verified-board, frontier, frontier+adapt, adapt-only,
frontier+flat, frontier+hact. Three closure rules: null (frontier
$\le$80\% of best-baseline cost-per-accepted), attribution (frontier+adapt
within 10\% of adapt-only), HACT (tree $\le$95\% of flat steel-man on
held-out evidence bytes).

\paragraph{Frozen outcomes (as generated 2026-09-23, never
recomputed).}
\begin{itemize}[leftmargin=*,itemsep=1pt]
\item \emph{Null closure triggered}: frontier cost-per-accepted
1{,}906{,}765 versus best baseline 564{,}773 (threshold $\le$451{,}818).
No byte-cost claim for frontier retirement.
\item \emph{Summary/reset falsified}: 21/63 accepted, 30{,}820 stale
reads, 280 lost obligations---the lossy-conversion failure mode.
\item \emph{Operations}: frontier cut retrieval operations 48.3\%
versus masking (27{,}728 vs 53{,}583 across the cohort), the
observation that became v2's regime-O hypothesis.
\item \emph{Attribution closure triggered}: scheduler dimension
byte-neutral (frontier+adapt identical to frontier on bytes); read
dimension-separated (makespan only; 87.4\% fewer rounds with adaptive
batching).
\item \emph{HACT closure not triggered}: held-out evidence bytes 59.8\%
lower than the flat steel-man pooled; per-family split
($+88.7\%$ clustered, $+49.2\%$ independent, $-87.1\%$ global) forced
the conditional reading formalized in v2's density rule.
\end{itemize}

\paragraph{Why v2 replaced the design.} v1's arms were hand-picked
combinations, so factor attribution leaned on pairwise comparisons and
a single attribution closure; its null rule read one pooled byte
threshold; its capsule accounting was an end-of-run total with no
independent audit; and it priced one cost regime, which is why the
operations result could only be reported as a hypothesis. Each weakness
became a v2 design item (Section~\ref{sec:retire-provenance}). v1's
cohort, closures, and summary remain reproducible byte-identically
from the committed v1 engine via
\texttt{bun src/harness/retirement/cli.ts}.


\bibliographystyle{plain}
\begin{thebibliography}{10}
\bibitem{delm} Yuzhen Mao and Azalia Mirhoseini.
\newblock Decentralized Multi-Agent Systems with Shared Context.
\newblock arXiv:2606.10662v1, June 2026.
\newblock \url{https://arxiv.org/abs/2606.10662v1}.

\bibitem{scaling} Yubin Kim et al.
\newblock Towards a Science of Scaling Agent Systems.
\newblock arXiv:2512.08296v3, revised April 8, 2026.
\newblock \url{https://arxiv.org/abs/2512.08296v3}.

\bibitem{context} Anthropic.
\newblock Effective Context Engineering for AI Agents.
\newblock Engineering article, September 29, 2025.
\newblock \url{https://www.anthropic.com/engineering/effective-context-engineering-for-ai-agents}.

\bibitem{alphabetic} Hirendu Vaishnav and Massoud Pedram.
\newblock Alphabetic Trees: Enumeration and Optimization with Applications to VLSI Design.
\newblock Author-hosted manuscript, Sections 4.1--4.4; publication date not established from the retrieved copy.
\newblock \url{https://www.mpedram.com/Papers/alp-trees-journal.pdf}.

\bibitem{baer} Michael B. Baer.
\newblock Alphabetic Coding with Exponential Costs.
\newblock arXiv:cs/0605099v3, revised March 2009.
\newblock \url{https://arxiv.org/abs/cs/0605099}.

\bibitem{hmt} Ziheng Shangguan, Aviv Yaish, and Dahlia Malkhi.
\newblock Authenticated Data Structures for Dynamic Workloads.
\newblock arXiv:2608.25206v1, August 2026.
\newblock \url{https://arxiv.org/abs/2608.25206v1}.

\bibitem{acar} Umut Acar, Guy Blelloch, Matthias Blume, Robert Harper, and Kanat Tangwongsan.
\newblock A Library for Self-Adjusting Computation.
\newblock \emph{Electronic Notes in Theoretical Computer Science}, 148:127--154, 2006.
\newblock \url{https://doi.org/10.1016/j.entcs.2005.11.043}.

\bibitem{pcaa} Zexun Wang.
\newblock Proof-Carrying Agent Actions: Model-Agnostic Runtime Governance for Heterogeneous Agent Systems.
\newblock arXiv:2606.04104, June 2026.
\newblock \url{https://arxiv.org/abs/2606.04104}.

\bibitem{hoeffding} Wassily Hoeffding.
\newblock Probability Inequalities for Sums of Bounded Random Variables.
\newblock \emph{Journal of the American Statistical Association}, 58(301):13--30, 1963.
\newblock \url{https://doi.org/10.1080/01621459.1963.10500830}.
\end{thebibliography}

\end{document}
